\documentclass[11pt,a4paper]{article}
\usepackage[T1]{fontenc}
\usepackage[utf8]{inputenc}
\usepackage{lmodern}
\usepackage{amsmath,amssymb,amsthm,mathtools}
\usepackage[a4paper,margin=24mm,headheight=15pt,footskip=27pt]{geometry}
\usepackage{microtype,booktabs,array,longtable,tabularx}
\usepackage{xcolor,enumitem,fancyhdr,titlesec,needspace}
\usepackage[most]{tcolorbox}
\usepackage{xurl,hyperref,bookmark,listings}
\definecolor{Ink}{HTML}{182B38}
\definecolor{Accent}{HTML}{125B70}
\definecolor{Pale}{HTML}{EFF5F7}
\definecolor{Rule}{HTML}{CCDCE2}
\hypersetup{colorlinks=true,linkcolor=Accent,citecolor=Accent,urlcolor=Accent,
 pdftitle={Six Sheets, Three Escapes: Complete Fiber Geometry and Explicit Normalization of a Five-Dimensional Keller Map},
 pdfauthor={Research prepared with ChatGPT for the ProveIt project (merged report of two AI-assisted manuscripts)},
 pdfsubject={Exact fibers, multiplicity patterns, normalization, discriminants and their local types, nonproperness, stabilization, and Galois groups of Gao's map F6}}
\urlstyle{same}
\titleformat{\section}{\Large\bfseries\color{Accent}}{\thesection}{0.6em}{}
\titleformat{\subsection}{\large\bfseries\color{Ink}}{\thesubsection}{0.6em}{}
\pagestyle{fancy}\fancyhf{}
\fancyhead[L]{\small Six Sheets, Three Escapes}
\fancyhead[R]{\small ProveIt research article}
\fancyfoot[C]{\thepage}
\renewcommand{\headrulewidth}{0.35pt}
\setlength{\parskip}{4pt plus 1pt}
\setlength{\parindent}{1.1em}
\setlength{\emergencystretch}{2em}
\setlist{itemsep=3pt,topsep=5pt}
\allowdisplaybreaks[2]
\newtheorem{theorem}{Theorem}[section]
\newtheorem{lemma}[theorem]{Lemma}
\newtheorem{proposition}[theorem]{Proposition}
\newtheorem{corollary}[theorem]{Corollary}
\theoremstyle{definition}
\newtheorem{definition}[theorem]{Definition}
\newtheorem{remark}[theorem]{Remark}
\newcommand{\A}{\mathbb A}
\newcommand{\C}{\mathbb C}
\newcommand{\Q}{\mathbb Q}
\newcommand{\R}{\mathbb R}
\newcommand{\Z}{\mathbb Z}
\newcommand{\Gm}{\mathbb G_m}
\newcommand{\Spec}{\operatorname{Spec}}
\newcommand{\Disc}{\operatorname{Disc}}
\newcommand{\Res}{\operatorname{Res}}
\newcommand{\Gal}{\operatorname{Gal}}
\newcommand{\Ram}{\operatorname{Ram}}
\newcommand{\simp}{\operatorname{simp}}
\newcommand{\codim}{\operatorname{codim}}
\newcommand{\rk}{\operatorname{rank}}
\newcommand{\id}{\operatorname{id}}
\newcommand{\F}{\mathcal F}
\newcommand{\base}{\mathcal A}
\newcommand{\normal}{\mathcal N}
\newcommand{\ord}{\mathcal O}
\newcommand{\kk}{\kappa}
\newcommand{\wt}{\widetilde}
\newcommand{\code}[1]{\texttt{\nolinkurl{#1}}}
\newcommand{\NP}{\operatorname{NP}}
\newcommand{\M}{M}
\newcommand{\dd}{\mathcal D}
\newcommand{\im}{\operatorname{im}}
\newcommand{\src}[1]{\textup{(source~#1)}}
\newtcolorbox{claimbox}{breakable,colback=Pale,colframe=Rule,boxrule=0.5pt,
 arc=1mm,left=2mm,right=2mm,top=1.5mm,bottom=1.5mm}
\lstset{basicstyle=\ttfamily\small,breaklines=true,columns=fullflexible,
 keepspaces=true,showstringspaces=false,frame=single,rulecolor=\color{Rule},
 backgroundcolor=\color{Pale},aboveskip=8pt,belowskip=8pt}
\begin{document}
\hypersetup{pageanchor=false}
\begin{titlepage}
\color{Ink}
\vspace*{12mm}
{\small\sffamily\bfseries\color{Accent} ALGEBRAIC GEOMETRY \; / \; GALOIS THEORY \; / \; EXACT CERTIFICATES\par}
\vspace{17mm}
{\Huge\bfseries Six Sheets,\par Three Escapes\par}
\vspace{8mm}
{\LARGE Complete Fiber Geometry and Explicit Normalization\par}
\vspace{3mm}
{\Large of a Five-Dimensional Keller Map\par}
\vspace{14mm}
{\large A research extension of the Jacobian-conjecture\par direction in the ProveIt repository\par}
\vspace{5mm}
{\normalsize Merged report of two manuscripts: \emph{Six Sheets, Three Escapes}
(source~02, the base) and \emph{A Codimension-Three Omitted Surface for a
Five-Dimensional Keller Map} (source~05)\par}
\vspace{5mm}
{Both sources dated 24 September 2026\par}
\vfill
\begin{claimbox}
\textbf{Principal result.} We determine every geometric fiber of Gao's map
$F_6$, including the previously untreated degeneration at its first target
coordinate zero. We construct its finite normalization with an explicit
six-element free basis, identify its omitted surface and its two-component
nonproperness locus, and prove geometric and arithmetic $S_6$ results.
Source~05 adds every root-multiplicity pattern, every local analytic type of
the discriminant, normal crossings with local monodromy $(\Z/2)^3$ along the
omitted surface, an obstruction to stabilization from dimension at most
three, and a second arithmetic $S_6$ certificate.

\textbf{Status.} The map $F_6$, its polynomiality and its generic degree are
Gao's prior work. The extensions are proved here with supplementary exact
computer-algebra checks. Both sources are AI-assisted research drafts
prepared with ChatGPT; neither is a peer-reviewed paper or a
Lean/Rocq-certified development, and nothing about $F_6$ is formalized in
ProveIt. Priority beyond the inspected sources has not been established.
\end{claimbox}
\vspace{6mm}
{\small Prepared for the ProveIt research program. Both sources' reproducibility
scripts, recorded outputs and sparse coefficient data accompany this report
(Appendix~\ref{f6:app:reproduce}).\par}
\end{titlepage}
\hypersetup{pageanchor=true}

\section*{Abstract}
\addcontentsline{toc}{section}{Abstract}
Gao's five-dimensional polynomial map $F_6$ has constant Jacobian determinant
$-290$ and generic geometric degree six. Its published construction provides
a sextic controlling fibers away from a degenerating target hyperplane,
but leaves the complete fiber stratification to further work. We resolve
the fiber-cardinality problem for this particular map and give additional
structural results.

The decisive identity is that the sweep parameter is the derivative of the
fiber sextic on its zero set. After a rescaling, this produces a uniform
root chart valid even on the exceptional hyperplane: the source open set
$x\ne0$ is exactly a simple-root locus. The remaining source hyperplane
maps isomorphically to the target hyperplane. Consequently the only
geometric fiber cardinalities are $0,1,2,3,4,6$, all of which occur. The
omitted set is an explicitly parametrized smooth closed surface isomorphic
to $\Gm\times\A^1$.

We then normalize the monogenic sextic order by adjoining one explicit
fraction. The resulting algebra is free of rank six, with three short
multiplication rules, and its spectrum is smooth. The original affine
space is the complement of the ramification locus and three unramified
boundary components in this finite model. Its nonproperness locus is
$V(cH)$, where $H$ is an irreducible normalized discriminant satisfying
$H|_{c=0}=-108(A^2+4B)$. A one-parameter Morse slice proves geometric
monodromy $S_6$; reductions modulo $3,13,37$ certify an explicit rational
specialization with the same Galois group. We also determine the possible
real fiber sizes, extract a general normalization lemma, and give an
exact inversion procedure and a proposed formalization and research agenda.

The second source, merged here, proves the fiber-size classification
independently and adds finer structure. Newton sums show that every sextic
of the family has at least three distinct roots, so exactly seven
root-multiplicity patterns can occur, and all seven do. The discriminant
surface in the three-dimensional coefficient space is normalized by the
affine plane, and its local analytic type is determined at every point. At
each point of the omitted surface the discriminant component of the
nonproperness locus has three smooth, mutually transverse branches, with
local monodromy $(\Z/2)^3$. The omitted surface excludes stabilization from
any polynomial self-map in at most three variables. The generic inverse
field extension has no proper intermediate field. A second Morse slice
reproves geometric $S_6$, and reductions modulo $7,11,269$ certify a second
rational specialization with Galois group $S_6$.

\medskip
\noindent\textbf{Keywords.} Keller map; polynomial fiber; normalization;
nonproperness; discriminant; monodromy; sextic; radical obstruction.

\newpage
\begingroup
\small
\setlength{\parskip}{0pt}
\tableofcontents
\endgroup
\newpage

\Needspace{6\baselineskip}
\section{The research target and the main conclusions}
\label{f6:sec:intro}

\Needspace{5\baselineskip}
\subsection{Relation to ProveIt and to the published problem}
The Jacobian-conjecture project in ProveIt contains explicit polynomial
maps, exact collision certificates, symmetry arguments, and stable
reductions of degree. Its research README studies degree and sparsity
questions for the three-dimensional example. The repository walkthrough
also points to Gao's account of the subsequent higher-dimensional
constructions \cite{ProveIt,Gao}. These give a natural opportunity to pass
from individual collision witnesses to the full geometry of an explicit
map.

We focus on the map denoted $F_6$ in \cite[\S4.5.1, Theorem 4.5]{Gao}.
The subscript labels the example; its \emph{ambient dimension is five}.
The map, its determinant, its sextic elimination equation, and the
four-point fibers over its nonzero distinguished axis are prior results.
Gao's concluding discussion and Appendix A.3 explicitly leave the
complete higher-dimensional stratifications, including that of $F_6$,
unresolved in version 1. We do not claim to have constructed a new
counterexample to the Jacobian conjecture. Our object is the full fiber
geometry and finite algebraic model of this existing example.

The source audit was scoped, not repository-wide. It covered the root
README, the Jacobian project and its research README, the relevant
walkthrough search result, and the cited primary paper. The exact
repository revision and inspected file hashes are recorded in
Appendix~\ref{f6:app:sources}. No claim is made that every research artifact
in this large repository has been searched for equivalent results.

\emph{Source~05's framing, kept.} Source~05 describes the same context: the
\code{Algebra/JacobianConjecture} development contains an explicit polynomial
counterexample, determinant and collision identities, and research on stable
reductions and equivariance, and its walkthrough points to Gao's broader
tangent-sweep construction as external mathematical context
\cite{ProveIt,ProveItWalkthrough}. (``Tangent-sweep'' is source~05's gloss;
the walkthrough does not use the word.) Source~05 states explicitly that
repository theorem names are not to be confused with theorems established
here. It also records what it does not claim: $F_6$ is not a newly
discovered construction; neither the plane Jacobian problem nor Gao's
separate middle-branch rigidity problem is settled; and no complete Whitney
stratification or irreducible decomposition of every higher-contact
incidence variety is claimed. Those are distinct, finer questions.

\Needspace{5\baselineskip}
\subsection{This merged report and ProveIt's formal development}
\label{f6:sec:merge}
This report is built from two manuscripts of ProveIt's incoming batch~36,
both pinned to ProveIt commit \code{e21766d04} and both dated
24~September~2026. Source~02, \emph{Six Sheets, Three Escapes}, is the base;
its text is kept, in its order, as Sections~\ref{f6:sec:intro}--\ref{f6:sec:conclusion}
and Appendices~\ref{f6:app:reproduce}--\ref{f6:app:sources}. Source~05,
\emph{A Codimension-Three Omitted Surface for a Five-Dimensional Keller Map},
proves the same spine (derivative identity, simple-root reconstruction,
spectrum, omitted surface, nonproperness, $S_6$) in different notation, and
adds results source~02 does not have. Every result, proof, example,
limitation and question of source~05 is printed here, in source~02's
notation (Section~\ref{f6:sec:dictionary}). A result proved by both is
printed once and credited to both; where source~05's proof is genuinely
different it is kept as a marked second route. Statements, proofs and
remarks taken from source~05 carry the tag \src{05} in their heading or
first sentence. Untagged text is source~02's, except the editorial text
written for this merge: Sections~\ref{f6:sec:merge} and~\ref{f6:sec:dictionary},
the rewritten Appendix~\ref{f6:app:reproduce}, Appendix~\ref{f6:app:provenance},
the re-scoping notes in Section~\ref{f6:sec:future}, and short linking
sentences that say which source proved what.
Appendix~\ref{f6:app:provenance} records the provenance and every choice the
merge made.

\emph{Formal status.} Nothing in this report is formalized. In particular,
the ProveIt Lean and Rocq developments contain no statement about $F_6$.
Their five-variable declarations, such as
\code{jacobianConjectureInDimensionFive_false}
(\code{Algebra/JacobianConjecture/Lean/JacobianConjecture/SimplerCounterexample.lean})
and \code{jacobian_conjecture_dimension_five_is_false}
(\code{Algebra/JacobianConjecture/Coq/SimplerCounterexample.v}), concern the
project's own stable map
$G=(G_1,\dots,G_5)$: a stable shear of Alpöge's three-dimensional map, of
ordinary degree six with profile $(6,6,4,3,4)$ and a two-point collision. That
$G$ is \emph{not} $F_6$. The map $F_6$ is Gao's \cite{Gao}; it has ambient
dimension five, geometric degree six, and ordinary coordinate degrees
$(7,38,40,42,44)$. This report belongs to ProveIt's research-report
collection and continues the formal project \code{Algebra/JacobianConjecture};
that relation confers no formal status on it. Source~05 says the same of its
own assertions: they have not been formalized in Lean or Rocq, have not
undergone external peer review, and do not inherit machine-checked status
from related ProveIt files; and its comparison with the pinned repository
and Gao's version~1 is not an exhaustive priority certification.

\Needspace{5\baselineskip}
\subsection{Conventions}
Unless indicated otherwise, $k$ is a field of characteristic zero. Equations
and chart isomorphisms are defined over $\Q$ and extend to $k$. A
\emph{geometric fiber} means the fiber after passage to an algebraic
closure of its residue field, and its cardinality counts distinct points,
not multiplicities. Statements about complex nonproperness and monodromy
use $k=\C$. Real and rational fibers are treated separately.

Write a source point as $(x,y,z_1,z_2,z_3)$ and a target point as
\[
 b=(c,A,B,D,E).
\]
The use of $D$ rather than $C$ for the fourth coordinate avoids confusion
with the field $\C$. Put
\begin{equation}\label{f6:eq:TL}
 \kk=E+B^2-AD,\qquad T=-A+2cB-c^3\kk,\qquad L=B-cD.
\end{equation}
The coordinate change from $E$ to $\kk$ is a polynomial automorphism of
target affine space.

Source~05 states its algebraic assertions about geometric fibers over an
algebraically closed field $k$ of characteristic zero and its analytic
assertions over $\C$; all its displayed formulas are defined over $\Q$. This
is the geometric-fiber setting above, so its proofs need no re-reading for
stronger hypotheses. It also fixes three words used throughout. A
\emph{Keller map} is a polynomial map with nonzero constant Jacobian
determinant. A \emph{simple root} has multiplicity exactly one; a repeated
root contributes \emph{no} point to the corresponding fiber, rather than one
point. The \emph{omitted locus} is the complement of the image; it can be
strictly smaller than the nonproperness locus (Definition before
Theorem~\ref{f6:thm:nonproper}).

\Needspace{8\baselineskip}
\subsection{Notation, and the dictionary between the two sources}
\label{f6:sec:dictionary}
Both sources use identical normalizations up to renaming; the merge checked
the dictionary below symbolically. This report uses source~02's names
throughout. Source~05's formulas are printed after the renaming in
Table~\ref{f6:tab:dictionary}; no normalization or scalar changes.

\begin{table}[htbp]
\centering\small
\renewcommand{\arraystretch}{1.25}
\begin{tabularx}{\textwidth}{@{}>{\raggedright\arraybackslash}p{0.27\textwidth}>{\raggedright\arraybackslash}p{0.3\textwidth}X@{}}
\toprule
Object & This report (= source 02) & Source 05\\
\midrule
target point & $b=(c,A,B,D,E)$ & $q=(C,p_1,p_2,p_3,p_4)$\\
scaled target & $Y_i=c^i\cdot(A,B,D,E)_i$ & $Y_i=C^ip_i$\\
stage monomials & $a_1,\dots,a_4$ & $\xi_1,\dots,\xi_4$\\
sweep variables & $(w_1,w_2,w_3)$ & $(t,s,r)$\\
potentials & $G_2,\ G_3,\ G_4$ & $A,\ 2tA+s,\ B$\\
sweep quotients $S_i/\gamma^i$ & $\widehat S_i(\gamma,u_1,u_2,u_3)$ & $E_i$\\
section coordinates at $x=0$ & $\beta_1,\dots,\beta_4$, right side of \eqref{f6:eq:section} & jets $J_1,\dots,J_4$\\
fiber sextic, root & $P_b(w)$, $w=\gamma u_1$ & $R_{a,b,c}(T)$, $T=\gamma u_1$\\
sextic coefficients & $(a,b,d)$ of \eqref{f6:eq:discsub} & $(a,b,c)$\\
universal discriminant & $\Delta(a,b,d)$ & $D(a,b,c)$\\
discriminant surface & $\dd=V(\Delta)\subset\A^3_{a,b,d}$ & $\mathcal D=V(D)$\\
normalized discriminant & $H=\Delta/c^2$ & $\widetilde D$\\
hyperplane quadratic's discriminant & $A^2+4B$ & $H:=p_1^2+4p_2$\\
free parameter on $M$ & $t=c^2B$ & $\eta=Y_2$\\
omitted surface & $M$ & $\mathcal M$\\
nonproperness set & $\NP(\F)$ & $S_{F}$\\
the cubic with $P=2Q_\star^2$ & $Q_\star(w)$ & $Q(T)$\\
triple-root cofactor; factors & $Q_\rho$;\ $h_4(\rho)$, $\Lambda(\rho)$ & $Q_\rho$;\ $H_4(\rho)$, $K(\rho)$\\
power sums of the roots & $N_1,N_2,N_3$ & $p_1,p_2,p_3$ (temporary)\\
two-root gap $\alpha-\beta$ & $\theta$ & $\delta$\\
local analytic coordinates & $(\sigma_1,\sigma_2,\sigma_3)$, quartic variable $\omega$ & $(u,v,w)$, $Z$\\
second Morse slice & $\hat p(w)+\tau$ & $P(T)+c$\\
second arithmetic target & $b_{**}=(1,-1,0,1,1)$ & $q_*$\\
\bottomrule
\end{tabularx}
\caption{Dictionary. Explicitly, $a=1+cA=Y_1+1$,
$b=c^4\kk-2c^2B+cA=Y_4+Y_2^2-Y_1Y_3-2Y_2+Y_1$ and
$d=c^3D-c^2B=Y_3-Y_2$.}
\label{f6:tab:dictionary}
\end{table}

\emph{Tempting false readings.} Several letters mean different things in the
two sources, and some also clash with ProveIt's project notation.
\begin{itemize}
\item $c$ is the \emph{first target coordinate} here; source~05's $c$ is the
constant coefficient of the sextic, here $d$. Likewise source~05's $b$, $a$
are coefficients, while here $b$ is also the target point (context decides).
\item $T$ here is the target function $-A+2cB-c^3\kk$ of \eqref{f6:eq:TL},
never a root variable; source~05's root variable $T$ is $w$ here.
\item $H$ here is the normalized discriminant (source~05's $\widetilde D$),
not source~05's $H=p_1^2+4p_2$, which is written $A^2+4B$ here.
\item $A,B$ here are target coordinates, not source~05's sweep potentials;
$D,E$ are target coordinates, not a discriminant and not quotient
polynomials. $E_1,E_2,E_3$ are the three escape components of
Theorem~\ref{f6:thm:boundary}.
\item $J(w)=g(w)+cA$ is a polynomial in $w$ here, not a boundary jet.
\item $\normal=\mathcal N$ denotes \emph{only} the normalization algebra of
$F_6$. Source~02 also wrote $\mathcal N_{\F}$ for the nonproperness set; this
report renames that set $\NP(\F)$ (merge disclosure).
\item ``Normalization'' has two objects: the normalization $Z=\Spec\normal$ of
the map $F_6$ (Section~\ref{f6:sec:normalization}), and the normalization
$\A^2_{a,\rho}\to\dd$ of the discriminant \emph{surface} in coefficient space
(Proposition~\ref{f6:kf:prop:Dnorm}). ``Normalized discriminant'' $H$ refers
to the first.
\item $t$ is the parameter on $M$ in \eqref{f6:eq:Mparam} and the Morse
parameter in Section~\ref{f6:sec:galois}; source~05's sweep parameter $t$ is
$w_1$ here.
\item Project notation: ProveIt's $F=(P,Q,R)$ is Alpöge's three-dimensional
map, and its $G=(G_1,\dots,G_5)$ is the five-variable stable map named above.
Neither is $F_6$; here $P_b$, $q_b$, $G_2,G_3,G_4$ are unrelated to them.
``Degree six'' for $F_6$ means generic fiber size six (geometric degree), not
the ordinary degree of the project's $G$.
\end{itemize}

\Needspace{5\baselineskip}
\subsection{A precise statement of what is proved}
Define
\begin{align}
 g(w)&=2w^4-2w^3-w+1=(w-1)(2w^3-1),\label{f6:eq:g}\\
 J(w)&=g(w)+cA,\label{f6:eq:J}\\
 P_b(w)&=w^2J(w)-cTw-c^2L,\label{f6:eq:P}\\
 q_b(v)&=v^2J(cv)-Tv-L.\label{f6:eq:qcompact}
\end{align}
Thus $P_b(cv)=c^2q_b(v)$. Expanded, the uniform polynomial is
\begin{align}
 q_b(v)={}&2c^4v^6-2c^3v^5-cv^3+(1+cA)v^2\notag\\
 &+(c^3\kk-2cB+A)v+cD-B.\label{f6:eq:q}
\end{align}

\begin{theorem}[Main conclusions]\label{f6:thm:main}
For Gao's map $\F=F_6$ the following hold.
\begin{enumerate}[label=\textup{(\roman*)}]
\item The source open set $x\ne0$ is isomorphic over target space to
\[
 \{(b,v):q_b(v)=0,\ q_b'(v)\ne0\}.
\]
The source hyperplane $x=0$ maps isomorphically onto $c=0$.
\item Every geometric fiber has
\[
 \#\F^{-1}(b)=\#\{\text{simple roots of }q_b\}+\mathbf1_{c=0}.
\]
For $c=0$ this equals $3$ when $A^2+4B\ne0$ and $1$ otherwise.
The complete set of attained cardinalities is $\{0,1,2,3,4,6\}$.
\item The omitted set over an algebraically closed field is the smooth
closed surface $M\simeq\Gm\times\A^1$ defined by
\begin{align}
 cA&=-\frac{11}{32},\notag\\
 c^3D-c^2B&=\frac{25}{128},\label{f6:eq:M}\\
 c^4\kk-2c^2B+cA&=\frac{5}{32}.\notag
\end{align}
In particular, $\F(\A^5)=\A^5\setminus M$ on geometric points.
\item The integral closure of the target coordinate ring in the source
function field is free of rank six with basis
\[
 1,w,w^2,w^3,w^4,\zeta.
\]
Its spectrum is smooth. Its multiplication is determined by
\begin{align}
 2w^5&=c\zeta+2w^4+w^2-(1+cA)w,\notag\\
 w\zeta&=Tw+cL,\label{f6:eq:rules}\\
 \zeta^2&=T\zeta+LJ(w).\notag
\end{align}
\item Let $H=\Disc_w(P_b)/c^2$. This is an irreducible polynomial, and
\[
 H(0,A,B,D,E)=-108(A^2+4B).
\]
The finite model's discriminant in the displayed basis is $H/256$.
The complex nonproperness locus of $\F$ is precisely $V(cH)$.
\item The geometric monodromy of the six-sheeted cover is $S_6$.
The rational target $(1,0,0,1,0)$ gives an irreducible sextic with Galois
group $S_6$ over $\Q$, so a full inverse vector there cannot be expressed
using radicals over $\Q$.
\end{enumerate}
\end{theorem}

The theorem will be proved in components. In particular, the displayed
sextic is \emph{not} itself the finite normalization: at $c=0$ it collapses
two branches to $w=0$. Ignoring this distinction gives wrong exceptional
fiber counts and an artificial discriminant component.

Source~05 states its main results as the following two theorems, printed
here in this report's notation. Parts~(ii) and~(iii) of
Theorem~\ref{f6:kf:thm:main} and the first two sentences of
Theorem~\ref{f6:kf:thm:geometry-summary} repeat parts~(ii), (iii), (v)
and~(vi) of Theorem~\ref{f6:thm:main}; both sources prove them.

\begin{theorem}[Complete fibers and image \src{05}]\label{f6:kf:thm:main}
For $\F=F_6$ the following hold.
\begin{enumerate}[label=\textup{(\roman*)}]
\item Every fiber is finite and reduced. If $c\ne0$, its points are in
explicit bijection with the simple roots of $P_b$. The only possible
root-multiplicity patterns of $P_b$, and the corresponding fiber sizes, are
\[
\begin{array}{c|ccccccc}
\text{pattern}&1^6&2\,1^4&3\,1^3&2^2\,1^2&4\,1^2&3\,2\,1&2^3\\ \hline
\#\F^{-1}(b)&6&4&3&2&2&1&0
\end{array}
\]
and each pattern occurs.
\item If $c=0$, then $\#\F^{-1}(b)=3$ when $A^2+4B\ne0$ and $1$ when
$A^2+4B=0$. The set of all fiber sizes is exactly $\{0,1,2,3,4,6\}$.
\item The complement of the image is exactly the surface $M$ of
\eqref{f6:eq:Mparam}. It is a smooth closed complete intersection in $\A^5$,
isomorphic to $\Gm\times\A^1$, of codimension three.
\item The map is not polynomial-equivalent to a stabilization of any
polynomial self-map in at most three variables.
\end{enumerate}
\end{theorem}

The geometry is stronger than a list of sample collisions: the theorem
identifies every target with no preimage, every possible loss of sheets, and
an exact algebraic decision rule for arbitrary targets.

\begin{theorem}[Escape and monodromy \src{05}]\label{f6:kf:thm:geometry-summary}
After the substitution \eqref{f6:eq:discsub}, $\Delta=c^2H$ with $H$
irreducible and $H(0,A,B,D,E)=-108(A^2+4B)$. Over $\C$ the nonproperness set
is exactly $V(cH)$, and the restriction of $\F$ over its complement is a
finite \'etale cover of degree six with monodromy group $S_6$. At each point
of $M$, the hypersurface $V(H)$ has three smooth, mutually transverse
analytic branches, and the local sheet exchanges form $(\Z/2\Z)^3$.
\end{theorem}

Where they are proved: part~(i) in Section~\ref{f6:kf:sec:patterns} and
Proposition~\ref{f6:kf:prop:polynomial}; part~(ii) in
Corollary~\ref{f6:cor:count} (with a second route in
Section~\ref{f6:kf:sec:hyperplane}); part~(iii) in
Theorem~\ref{f6:thm:omitted}; part~(iv) in
Corollary~\ref{f6:kf:cor:nostabilization}; the escape statements in
Section~\ref{f6:sec:discriminant}, including Corollary~\ref{f6:kf:cor:crossings};
and monodromy in Theorem~\ref{f6:thm:geometricS6}.

\Needspace{6\baselineskip}
\section{The polynomial map, without its large expansion}
\label{f6:sec:map}

\Needspace{5\baselineskip}
\subsection{A compact definition}
We reproduce the construction needed to identify the map unambiguously
\cite[\S4.5.1]{Gao}. Set
\[
 a_1=xy,\quad a_2=x^2z_1,\quad a_3=x^3z_2,\quad a_4=x^4z_3,
\]
and apply the polynomial stage
\begin{align}
 \gamma&=1-29a_1+999a_1^2+355a_1a_2-41553a_1^3+a_4,\notag\\
 u_1&=1+27a_1-5a_2,\label{f6:eq:stage}\\
 u_2&=a_1-12a_2+\frac{2128}{5}a_1^2+a_3,\notag\\
 u_3&=-4a_1-10a_2.\notag
\end{align}
Introduce sweep variables
\[
 w_1=\gamma u_1,\qquad w_2=\gamma^3u_2,\qquad w_3=\gamma^4u_3.
\]
Define potentials and one scalar polynomial by
\begin{align}
 G_2&=w_2+w_1w_3+w_1^2-w_1^3,\notag\\
 G_3&=2w_1G_2+w_2,\notag\\
 G_4&=w_3-G_2^2+2w_1^4-2w_1^5,\label{f6:eq:potentials}\\
 X_1&=-w_3+2w_2-2w_1+2w_1w_3+5w_1^2-2w_1^3
                +8w_1^4-10w_1^5.\notag
\end{align}
Put
\begin{align*}
 S_1&=X_1+\gamma,& S_2&=G_2+w_1S_1,\\
 S_3&=G_3+w_1^2S_1,& S_4&=G_4+w_2S_1.
\end{align*}
With $c=x\gamma$, the map is
\begin{equation}\label{f6:eq:map}
 \F(x,y,z_1,z_2,z_3)=
 \left(c,\frac{S_1}{c},\frac{S_2}{c^2},\frac{S_3}{c^3},\frac{S_4}{c^4}\right).
\end{equation}
All quotients in \eqref{f6:eq:map} cancel \emph{as polynomials}. This
cancellation is part of the construction, not a restriction to $c\ne0$.

\begin{proposition}[Prior polynomiality, with an independent certificate]
\label{f6:prop:keller}
The expression \eqref{f6:eq:map} defines a polynomial map over $\Q$, of
coordinate degrees $(7,38,40,42,44)$, and
$\det J_{\F}=-290$. The expanded coordinates contain respectively
$6,342,421,507,904$ nonzero monomials.
\end{proposition}
\begin{proof}
Polynomiality and the degree profile are established by the stage and
weighted-divisibility construction in \cite[Theorem 4.5]{Gao}.
For an independent finite check, the accompanying script first substitutes
$(w_1,w_2,w_3)=(\gamma u_1,\gamma^3u_2,\gamma^4u_3)$ into each $S_i$
and divides each monomial by $\gamma^i$. It then substitutes
\eqref{f6:eq:stage}, verifies that each resulting monomial has $x$-exponent
at least $i$, and divides by $x^i$. Thus no rational evaluation or
numerical cancellation is used. Counting the resulting sparse terms gives
the stated data.

For the determinant, the monomial change to the $a_i$ has determinant
$x^{10}$. The stage has determinant $-290$; its inverse is displayed
below. The change to the $w_i$ has determinant $\gamma^8$. Direct
four-variable differentiation gives
\[
 \det\frac{\partial(S_1,S_2,S_3,S_4)}
               {\partial(\gamma,w_1,w_2,w_3)}=\gamma.
\]
Padding the sweep with $c=x\gamma$ contributes another $\gamma$.
Finally the target rescaling has determinant $c^{-10}$. On
$x\gamma\ne0$ their product is
\[
 x^{10}(-290)\gamma^8\gamma^2c^{-10}=-290.
\]
Both sides are polynomials, so the identity holds everywhere. The small
stage and sweep determinants are also independently checked in the
certificate.
\end{proof}

\begin{proof}[Second proof, by finite jets \src{05}]
\label{f6:kf:prop:polynomial}
Source~05 states this as a proposition: after the stage substitution, $x^i$
divides $\widehat S_i$ for $i=1,2,3,4$, and the resulting polynomial map
\[
 \F=\Bigl(x\gamma,\frac{\widehat S_1}{x},\frac{\widehat S_2}{x^2},
 \frac{\widehat S_3}{x^3},\frac{\widehat S_4}{x^4}\Bigr)
\]
has $\det J_\F=-290$. (This is \eqref{f6:eq:map}, since $S_i=\gamma^i\widehat
S_i$ and $c=x\gamma$.) Here $\widehat S_i\in\Q[\gamma,u_1,u_2,u_3]$ is the
quotient of $S_i$ by $\gamma^i$ after the substitution
$(w_1,w_2,w_3)=(\gamma u_1,\gamma^3u_2,\gamma^4u_3)$; the four quotients are
short, and are printed in Appendix~\ref{f6:kf:app:polynomials}. Two useful
identities, obtained by differentiating the displayed polynomials, are
\begin{equation}\label{f6:kf:eq:sweepdet}
 X_1=-\det\frac{\partial(G_2,G_3,G_4)}{\partial(w_1,w_2,w_3)},
 \qquad
 \det\frac{\partial(S_1,S_2,S_3,S_4)}{\partial(\gamma,w_1,w_2,w_3)}=\gamma .
\end{equation}

For divisibility by powers of $x$, substitute
\begin{align*}
\gamma&=1-29xy+999x^2y^2+x^3(355yz_1-41553y^3)+x^4z_3,\\
u_1&=1+27xy-5x^2z_1,\\
u_2&=xy+x^2\Bigl(-12z_1+\frac{2128}{5}y^2\Bigr)+x^3z_2,\\
u_3&=-4xy-10x^2z_1
\end{align*}
into the appendix polynomials, and collect only the coefficients through
$x^4$. This gives
\begin{equation}\label{f6:kf:eq:jets}
 \widehat S_i=x^i\beta_i+O(x^{i+1}),
\end{equation}
where $(\beta_1,\beta_2,\beta_3,\beta_4)$ are exactly the right-hand sides of
\eqref{f6:eq:section} for $A,B,D,E$. (Source~05 prints the coefficient of
$y^2z_1$ in $\beta_4$ as $-463585/25$, which is source~02's $-92717/5$.) Here
$O(x^m)$ means membership in the polynomial ideal $(x^m)$, not an analytic
asymptotic assertion. All coefficients below $x^i$ are zero, so this is a
finite coefficient proof of divisibility; it can be performed in
$\Q[y,z_1,z_2,z_3][x]/(x^5)$, where terms of degree at least five cannot
affect any of these claims. That is why a truncated multiplication is exact
rather than heuristic.

On the nonempty open set $x\gamma\ne0$, the map is a composition of five
maps with Jacobian factors, in order,
\[
 x^{10},\qquad -290,\qquad \gamma^8,\qquad\gamma^2,\qquad (x\gamma)^{-10}.
\]
The first comes from the monomials $a_i$, the second from the stage, the third
from the scaling of the sweep variables. For the fourth, the padded sweep
sends $(x,\gamma,w_1,w_2,w_3)$ to $(x\gamma,S)$; expansion in the $x$ column
gives $\gamma\det J S=\gamma^2$. The last map divides $S_i$ by
$(x\gamma)^i$, keeping $x\gamma$ fixed. The product is $-290$; both sides
are polynomials, and equality on a dense open set proves equality
everywhere.
\end{proof}

\begin{remark}[Every fiber point is reduced \src{05}]\label{f6:kf:rem:etale}
In particular $\F$ is \'etale: its differential is invertible everywhere.
Once finiteness of its fibers is established, every fiber point has
multiplicity one. More concretely, at a fiber point the linear terms of the
five defining equations generate the cotangent space, so the
zero-dimensional local fiber ring is the residue field.
\end{remark}

\Needspace{5\baselineskip}
\subsection{The stage inverse and the source hyperplane}
The stage inverse is particularly useful. Given
$(\gamma,u_1,u_2,u_3)$, recover
\begin{align}
 a_1&=\frac{2u_1-2-u_3}{58},&
 a_2&=\frac{27a_1-u_1+1}{5},\notag\\
 a_3&=u_2-a_1+12a_2-\frac{2128}{5}a_1^2,&
 a_4&=\gamma-1+29a_1-999a_1^2-355a_1a_2+41553a_1^3.
 \label{f6:eq:stageinverse}
\end{align}
Substitution proves both compositions are the identity.
Source~05 writes the second entry as $a_2=-\bigl(4(u_1-1)+27u_3\bigr)/290$,
which is the same function: the first two equations solve a $2\times2$ linear
system of determinant $-290$, and the last two are triangular substitutions.
Direct differentiation, in the orders $(a_1,a_2,a_3,a_4)$ and
$(\gamma,u_1,u_2,u_3)$, gives stage determinant $-290$ \src{05}.

The polynomial restriction of $\F$ to $x=0$ is
\begin{align}
 c&=0,\notag\\
 A&=y,\notag\\
 B&=-\frac{4636}{5}y^2+29z_1,\label{f6:eq:section}\\
 D&=\frac{6039}{5}y^3-1041yz_1+5z_2,\notag\\
 E&=\frac{74000249}{25}y^4-\frac{92717}{5}y^2z_1
       +5yz_2-1041z_1^2-2z_3.\notag
\end{align}
This identity can be obtained by taking the coefficients of $x^i$ before
the final divisions; the accompanying polynomial construction checks all
five coordinates independently. The diagonal coefficients $1,29,5,-2$
are nonzero, so \eqref{f6:eq:section} is a triangular polynomial isomorphism.
It provides exactly one source point on $x=0$ over every target on $c=0$.
Its determinant $-290$ also checks the global determinant normalization.

\Needspace{6\baselineskip}
\section{The derivative identity and the uniform inverse chart}
\label{f6:sec:chart}

\Needspace{5\baselineskip}
\subsection{Eliminating the sweep variables}
For a target $b$ set
\[
 Y_1=cA,\qquad Y_2=c^2B,\qquad Y_3=c^3D,\qquad Y_4=c^4E.
\]
Write $w=w_1$. The equations $S_i=Y_i$ and
\eqref{f6:eq:potentials} give
\begin{align}
 w_2&=Y_3-2wY_2+w^2Y_1=:V_Y(w),\label{f6:eq:V}\\
 w_3&=Y_4+Y_2^2-Y_1Y_3+2w^5-2w^4=:W_Y(w).\label{f6:eq:W}
\end{align}
Indeed $G_2=Y_2-wY_1$, and eliminating $G_3$ first gives
\eqref{f6:eq:V}; eliminating $G_4$ next gives \eqref{f6:eq:W}.
Substitution into $G_2=w_2+ww_3+w^2-w^3$ leaves
\begin{align}
 R_Y(w)={}&2w^6-2w^5-w^3+(Y_1+1)w^2\notag\\
 &+(Y_4+Y_2^2-Y_1Y_3-2Y_2+Y_1)w+Y_3-Y_2=0.
 \label{f6:eq:RY}
\end{align}
This is Gao's sextic. Substituting the scaled target gives exactly $P_b$.

\begin{lemma}[Derivative identity]\label{f6:lem:derivative}
As a polynomial identity in $Y$ and $w$,
\begin{equation}\label{f6:eq:derivative}
 Y_1-X_1\bigl(w,V_Y(w),W_Y(w)\bigr)=R_Y'(w)-2R_Y(w).
\end{equation}
Consequently, on the fiber equation, $\gamma=R_Y'(w)$.
\end{lemma}
\begin{proof}
Insert \eqref{f6:eq:V} and \eqref{f6:eq:W} into the displayed expression for
$X_1$ in \eqref{f6:eq:potentials} and collect coefficients in $w$. The result
is $R_Y'-2R_Y$. This is an identity before imposing $R_Y=0$, so it
also specifies the exact ideal membership, not just agreement at selected
roots. Finally $S_1=Y_1$ says $\gamma=Y_1-X_1$.
\end{proof}

\begin{proof}[Second proof, with the expansion displayed \src{05}]
Source~05 proves the identical lemma by printing the collected polynomial.
Let $B_0=Y_4+Y_2^2-Y_1Y_3$, so the linear coefficient of $R_Y$ is
$B_0-2Y_2+Y_1$. Substitution into $X_1$ gives
\begin{align*}
X_1\bigl(w,V_Y(w),W_Y(w)\bigr)={}&4w^6-16w^5+10w^4-2w^3+(5+2Y_1)w^2\\
&+(2B_0-4Y_2-2)w-B_0+2Y_3.
\end{align*}
Subtract this from $Y_1$ and compare with
$R_Y'=12w^5-10w^4-3w^2+2(Y_1+1)w+B_0-2Y_2+Y_1$. The difference is exactly
$-2R_Y$, coefficient by coefficient. In the elimination no degree drops
occur: the leading coefficient of $R_Y$ is the constant~$2$.
\end{proof}

Source~05 calls this the main simplification: an arbitrary repeated root is
deleted by the same condition, not merely by an extra filter checked at
selected targets.

For $c\ne0$, the relation $c=x\gamma$ forces $\gamma\ne0$.
Thus exactly the \emph{simple} sextic roots survive. A multiple root is
not a multiple point of the original fiber: it produces no source point.
The distinction matters because $\F$ itself is everywhere \'etale.

\Needspace{5\baselineskip}
\subsection{A rescaling that retains the exceptional fibers}
On $x\ne0$ use the regular source function
\begin{equation}\label{f6:eq:vforward}
 v=\frac{u_1}{x}=\frac1x+27y-5xz_1.
\end{equation}
When $c\ne0$ this is $w/c$. The identities $P_b(cv)=c^2q_b(v)$ and
$P_b'(cv)=cq_b'(v)$ suggest that the derivative of $q_b$, rather than
a quotient by $c$, is the correct denominator for inversion.

\begin{theorem}[Uniform simple-root chart]\label{f6:thm:chart}
Let
\[
 \mathcal U=\Spec k[c,A,B,D,E,v,1/q_b'(v)]/(q_b(v)).
\]
There is an isomorphism over target space
\[
 \A^5_{x\ne0}\ \simeq\ \mathcal U.
\]
In one direction use \eqref{f6:eq:vforward}. Conversely, for a simple root
$v$ set $\delta=q_b'(v)$ and define
\begin{align}
 x&=\delta^{-1},& \gamma&=c\delta,\notag\\
 u_1&=v/\delta,&
 u_2&=(D-2vB+v^2A)/\delta^3,\label{f6:eq:inversechart}\\
 u_3&=(\kk+2cv^5-2v^4)/\delta^4.&&\notag
\end{align}
Recover $a_i$ using \eqref{f6:eq:stageinverse}, and then put
\begin{equation}\label{f6:eq:sourceinverse}
 y=a_1\delta,\quad z_1=a_2\delta^2,\quad
 z_2=a_3\delta^3,\quad z_3=a_4\delta^4.
\end{equation}
No division by $c$ occurs.
\end{theorem}
\begin{proof}
First work on $c\ne0$. Lemma~\ref{f6:lem:derivative} gives
$\gamma=P_b'(cv)=c\delta$, hence $x=1/\delta$.
Equations \eqref{f6:eq:V}--\eqref{f6:eq:W} give
\[
 w_2=c^3(D-2vB+v^2A),\qquad
 w_3=c^4(\kk+2cv^5-2v^4).
\]
Dividing by $\gamma^3$ and $\gamma^4$, respectively, gives
\eqref{f6:eq:inversechart}. The stage inverse and monomial inverse now give
\eqref{f6:eq:sourceinverse}. These steps reverse the elimination, so they
prove both compositions on $c\ne0$.

We justify extension across $c=0$, rather than assuming it. The polynomial
$q_b$ is irreducible in $k[c,A,B,D,E,v]$: as a polynomial in $E$ it is
linear with coefficient $c^3v$, and that coefficient is coprime to its
constant term. To see the coprimality, specialize $c=0$ and $v=0$:
\[
 q_b|_{c=0}=v^2+Av-B,\qquad q_b|_{v=0}=cD-B.
\]
Neither $c$ nor $v$ divides the constant term. Gauss's lemma applied to
the fraction field of the other variables proves irreducibility. Thus
$\mathcal U$ is integral and its $c\ne0$ locus is dense.

On the source open set $x\ne0$, the functions in \eqref{f6:eq:vforward}
and $\F$ are regular. The relations $q_b(v)=0$ and
$q_b'(v)=1/x$ hold on the dense open set $c\ne0$, so they hold throughout.
Consequently the forward map lands in $\mathcal U$. All inverse formulas
are regular on $\mathcal U$, and the two composition identities extend
from its dense $c\ne0$ locus. This proves the isomorphism everywhere.
\end{proof}

\begin{theorem}[Uniform reconstruction over $c\ne0$ \src{05}]
\label{f6:kf:thm:reconstruction}
For $c\ne0$, the fiber over $b$ is naturally isomorphic to
\begin{equation}\label{f6:kf:eq:fiberalgebra}
 \Spec k[w,1/P_b'(w)]/(P_b(w)).
\end{equation}
The bijection on points sends a simple root $w$ to the source point
\begin{align}
\gamma&=P_b'(w),& x&=c/\gamma,\notag\\
u_1&=w/\gamma,& u_2&=V_Y(w)/\gamma^3,& u_3&=W_Y(w)/\gamma^4;
\label{f6:kf:eq:reconstruction}
\end{align}
then \eqref{f6:eq:stageinverse} gives $a_1,\dots,a_4$, and
$y=a_1/x$, $z_1=a_2/x^2$, $z_2=a_3/x^3$, $z_3=a_4/x^4$. These formulas
work in families over the entire open set $c\ne0$.
\end{theorem}
\begin{proof}[Proof \src{05}]
A source point over such a target has $x\gamma\ne0$. Its sweep parameter
$w=\gamma u_1$ satisfies $P_b(w)=0$ by the elimination, and
$P_b'(w)=\gamma\ne0$ by Lemma~\ref{f6:lem:derivative}. Conversely, a simple
root makes every denominator in \eqref{f6:kf:eq:reconstruction} nonzero. The
equation $P_b=0$ recovers $G_2=Y_2-wY_1$; the definitions of $V_Y$, $W_Y$
recover the other two potential equations; identity
\eqref{f6:eq:derivative} recovers the first sweep equation. The inverse
stage and monomial substitutions give a point of the required fiber. Each
operation inverts the corresponding forward operation, so this is a
bijection. The same identities hold in coordinate rings after inverting $c$
and $P_b'$, so they prove the scheme isomorphism
\eqref{f6:kf:eq:fiberalgebra}, not merely a set-theoretic statement for
closed points. For a root of multiplicity at least two, $P_b'$ belongs to
the maximal ideal of its local Artinian factor, and localization at $P_b'$
annihilates that entire factor; a simple-root factor survives as $k$.
\end{proof}

This is the $c\ne0$ part of Theorem~\ref{f6:thm:chart}, with $w=cv$ and
$x=c/P_b'(cv)=1/q_b'(v)$. Source~05 treats $c=0$ pointwise instead
(Section~\ref{f6:kf:sec:hyperplane}); source~02's chart is the stronger,
uniform statement.

\Needspace{5\baselineskip}
\subsection{The three effective parameters and a finite completion
over \texorpdfstring{$c\ne0$}{c nonzero} \src{05}}
\label{f6:kf:sec:parameters}
On $c\ne0$ the change from $Y$ to $(a,b,d,t)$, with $t=Y_2$, is a polynomial
automorphism. Its inverse is
\begin{align}
Y_1&=a-1,\qquad Y_2=t,\qquad Y_3=d+t,\notag\\
Y_4&=b-t^2+(a-1)(d+t)+2t-(a-1).\label{f6:kf:eq:Yinv}
\end{align}
Hence on $c\ne0$ the base is $\Gm\times\A^1_t\times\A^3_{a,b,d}$, and the
fiber polynomial depends only on $(a,b,d)$. Every parameter triple is
realized by actual targets, for example with $c=1$ and $t=0$. (This $t$ is
the parameter of the omitted surface in \eqref{f6:eq:Mparam}.)

There is also a finite completion over this open base:
\begin{equation}\label{f6:kf:eq:completion}
 \mathcal X^*=\Spec k[c,c^{-1},t,a,b,d,w]/(P),\qquad
 P=2w^6-2w^5-w^3+aw^2+bw+d .
\end{equation}
It is finite free of rank six over the base, with basis $1,w,\dots,w^5$.
Since $\partial P/\partial d=1$, it is smooth and isomorphic to
$\Gm\times\A^4$ after eliminating $d$. The actual source over $c\ne0$ is
exactly the open subset $P'\ne0$ of this completion. Thus repeated-root
points do exist in the finite completion: they are precisely the points
excluded from the affine source by reconstruction. Over $c\ne0$ this is the
finite model $Z$ of Section~\ref{f6:sec:normalization}, since
$\normal[1/c]=\ord[1/c]$; source~05 asked for its extension across $c=0$,
which Theorems~\ref{f6:thm:normalization} and~\ref{f6:thm:boundary} supply
(Section~\ref{f6:kf:sec:future}, direction~11).

\begin{corollary}[All geometric fibers]\label{f6:cor:count}
For every target,
\begin{equation}\label{f6:eq:count}
 \#\F^{-1}(b)=\#\{v:q_b(v)=0,\ q_b'(v)\ne0\}+\mathbf1_{c=0}.
\end{equation}
For $c=0$,
\begin{equation}\label{f6:eq:c0count}
 \#\F^{-1}(0,A,B,D,E)=
 \begin{cases}3,&A^2+4B\ne0,\\1,&A^2+4B=0.\end{cases}
\end{equation}
\end{corollary}
\begin{proof}
Partition the source into $x\ne0$ and $x=0$. Apply
Theorem~\ref{f6:thm:chart} to the first part and \eqref{f6:eq:section} to the
second. At $c=0$, the polynomial is $v^2+Av-B$, with derivative $2v+A$.
A nonzero discriminant gives two simple roots. A zero discriminant gives
one double root, which is deleted; it contributes no point.
\end{proof}

\Needspace{5\baselineskip}
\subsection{The exceptional hyperplane directly on the source \src{05}}
\label{f6:kf:sec:hyperplane}
Source~05 proves \eqref{f6:eq:c0count} without the rescaled chart, by
splitting the source according to $x\gamma=0$. The open-set calculation must
not be specialized blindly to $c=0$: there the sextic becomes
$w^2g(w)$, independent of the target coordinates, whereas the actual fiber
size is not constant.

\emph{The unique point with $x=0$.} At $x=0$ the stage gives $\gamma=1$, and
the jets \eqref{f6:kf:eq:jets} give the section \eqref{f6:eq:section}, a
triangular automorphism of the hyperplane with diagonal $1,29,5,-2$. Its
inverse is found in the order
\begin{align*}
y&=A,\qquad z_1=\frac{B+(4636/5)y^2}{29},\qquad
z_2=\frac{D+1041yz_1-(6039/5)y^3}{5},\\
z_3&=-\frac12\Bigl(E-5yz_2+1041z_1^2+\frac{92717}{5}y^2z_1
     -\frac{74000249}{25}y^4\Bigr).
\end{align*}
In particular, the origin has one explicitly visible preimage, the origin.

\emph{The branch $\gamma=0$, $x\ne0$.} From the quotient polynomials of
Appendix~\ref{f6:kf:app:polynomials},
\begin{equation}\label{f6:kf:eq:Egamma0}
(\widehat S_1,\widehat S_2,\widehat S_3,\widehat S_4)|_{\gamma=0}
=\bigl(1-2u_1,\ u_1-u_1^2,\ u_1^2+u_2,\ u_1^4+(1-2u_1)u_2+u_3\bigr).
\end{equation}
The first two target equations give
\begin{equation}\label{f6:kf:eq:xboundary}
u_1=\frac{1-xA}{2},\qquad x^2(A^2+4B)=1,
\end{equation}
and the last two give
\begin{equation}\label{f6:kf:eq:u23boundary}
u_2=x^3D-u_1^2,\qquad u_3=x^4E-(1-2u_1)u_2-u_1^4.
\end{equation}
For each solution $x\ne0$ of \eqref{f6:kf:eq:xboundary}, set $\gamma=0$ and
use the stage inverse \eqref{f6:eq:stageinverse}; then recover
$y,z_1,z_2,z_3$ by dividing the $a_i$ by the appropriate powers of $x$. This
supplies a unique point and proves sufficiency of every displayed equation.
If $A^2+4B\ne0$, the equation $x^2(A^2+4B)=1$ has exactly two distinct
solutions over $k$; if $A^2+4B=0$, it has none. Together with the single
$x=0$ point this proves \eqref{f6:eq:c0count}. The two branches do not
overlap, since $\gamma=1$ on $x=0$. (In the chart of
Theorem~\ref{f6:thm:chart} these two points are the simple roots $v=u_1/x=(1-xA)/(2x)$ of
$v^2+Av-B$.)

Every fiber is therefore finite, with the claimed cardinality, and
Remark~\ref{f6:kf:rem:etale} gives reducedness. With
Theorem~\ref{f6:kf:thm:reconstruction} and Section~\ref{f6:kf:sec:patterns},
this is source~05's proof of parts~(i) and~(ii) of
Theorem~\ref{f6:kf:thm:main}.

\Needspace{6\baselineskip}
\section{Every fiber cardinality and the omitted surface}
\label{f6:sec:fibers}

\Needspace{5\baselineskip}
\subsection{A finite algebraic test at every target}
For a nonzero characteristic-zero polynomial $f$ define its simple-root
factor by
\begin{equation}\label{f6:eq:simplepart}
 G=\gcd(f,f'),\qquad S=f/G,\qquad
 \simp(f)=\frac{S}{\gcd(S,G)},
\end{equation}
with gcds and final polynomials normalized to be monic. A root of
multiplicity $m$ occurs in $G$ with multiplicity $m-1$, and in $S$ with
multiplicity one. The second gcd removes it precisely when $m\ge2$.
Hence
\[
 \#\{\text{simple geometric roots of }f\}=\deg\simp(f).
\]
Together with \eqref{f6:eq:count}, this is an exact complete fiber-count
criterion, including singular discriminant strata. It does not require
factorization over an algebraic closure, and it is not limited to generic
points of a discriminant component.

Source~05 gives the same test independently \src{05}: with $f=P_b/2$ monic
and all gcds monic, put $d_1=\gcd(f,f')$, $h=f/d_1$ and
$s_1=h/\gcd(h,d_1)$; then $\#\F^{-1}(b)=\deg s_1$ for $c\ne0$. Together with
\eqref{f6:eq:c0count} this is a terminating exact algorithm for every rational
or effectively represented algebraic target. It counts points over an
algebraic closure and does not assert that those points are rational. It uses
no numerical root-separation tolerance and no multivariate Gröbner-basis
computation.

\begin{theorem}[Cardinality spectrum]\label{f6:thm:spectrum}
The geometric fiber cardinalities of $\F$ are exactly
\[
 \boxed{\{0,1,2,3,4,6\}.}
\]
\end{theorem}
\begin{proof}
For $c\ne0$, $P_b$ has degree six. Exactly five simple roots are
impossible: the remaining degree one would have to be another simple
root. Thus its number of simple roots belongs to
$\{0,1,2,3,4,6\}$. For $c=0$, \eqref{f6:eq:c0count} gives only $1$ and $3$.
Table~\ref{f6:tab:witness} supplies exact witnesses for every listed value.
For nonzero $c$ in the table, $c=1$, so $P_b=q_b$ with the variable
renamed. The displayed factorizations and elementary gcds give the
counts; the six-point row has discriminant $-993580\ne0$.
\end{proof}

\begin{corollary}[The missing five-point fiber \src{05}]\label{f6:kf:cor:no5}
Over $c\ne0$, a fiber has either six points or at most four points. A
repeated root removes at least two of the six roots, and cannot remove just
one.
\end{corollary}

\begin{table}[htbp]
\centering\small
\renewcommand{\arraystretch}{1.4}
\begin{tabularx}{\textwidth}{@{}c>{\raggedright\arraybackslash}p{0.32\textwidth}X@{}}
\toprule
Size & Target $(c,A,B,D,E)$ & Exact fiber polynomial or criterion\\
\midrule
$0$ & $(1,-11/32,0,25/128,1773/4096)$ &
$2(w^3-w^2/2-w/8-5/16)^2$\\
$1$ & $(0,0,0,0,0)$ & $q(v)=v^2$; only the section remains\\
$2$ & $(1,-11/8,0,-1/16,171/128)$ &
$(w^2+1/4)^2(2w^2-2w-1)$\\
$3$ & $(1,-8,0,-7,79)$ &
$(w-1)^3(2w^3+4w^2+6w+7)$\\
$4$ & $(1,0,0,0,0)$ & $w^2(w-1)(2w^3-1)$\\
$6$ & $(1,0,0,1,0)$ & $2w^6-2w^5-w^3+w^2+1$\\
\bottomrule
\end{tabularx}
\caption{Exact witnesses for the entire geometric cardinality spectrum.}
\label{f6:tab:witness}
\end{table}

\Needspace{5\baselineskip}
\subsection{Classifying all sextics with no surviving root}
The empty-fiber case can be solved in closed form, not merely by the gcd
criterion. Divide $P_b$ by $2$. Its three leading nontrivial coefficient
constraints are
\begin{equation}\label{f6:eq:fixedcoeff}
 \frac{P_b(w)}2=w^6-w^5+0w^4-\frac12w^3+\text{terms of degree at most two}.
\end{equation}
A sextic with no simple root has multiplicity partition
$6$, $4+2$, $3+3$, or $2+2+2$.

\begin{lemma}[Unique all-multiple member]\label{f6:lem:square}
A polynomial of the form \eqref{f6:eq:fixedcoeff} has no simple root if and
only if
\begin{equation}\label{f6:eq:uniquesquare}
 P_b(w)=2\left(w^3-\frac12w^2-\frac18w-\frac5{16}\right)^2.
\end{equation}
The cubic in parentheses has three distinct roots.
\end{lemma}
\begin{proof}
For multiplicity $6$, write $(w-r)^6$. The $w^5$ coefficient forces
$r=1/6$, whereas the $w^4$ coefficient is then $15r^2\ne0$.

For multiplicity $4+2$, write $(w-r)^4(w-s)^2$. The $w^5$ coefficient
forces $s=1/2-2r$. The next two required equations are
\[
 -\frac{24r^2-8r-1}{4}=0,
 \qquad \frac{8r^3+4r^2-2r+1}{2}=0.
\]
The remainder of the second polynomial on division by the first is
$(10r+23)/36$. Thus $r=-23/10$, but
$24r^2-8r-1=3609/25\ne0$, a contradiction.

For multiplicity $3+3$, write $(w^2-sw+p)^3$. Matching $w^5$ and $w^4$
gives $s=1/3$ and $p=-1/9$. Its $w^3$ coefficient is then $5/27$, not
$-1/2$.

The remaining case is a square of a monic cubic. Write
$(w^3+aw^2+bw+d)^2$ and compare $w^5,w^4,w^3$ in that order. This gives
$a=-1/2$, $b=-1/8$, $d=-5/16$, proving uniqueness. Conversely its square
has no simple root. The cubic discriminant is $-401/128\ne0$, so the
three roots are distinct and all have multiplicity exactly two in $P_b$.
\end{proof}

Source~05 excludes more than the all-multiple partitions: it excludes every
pattern with fewer than three distinct roots, and then reproves the lemma.
In its words, the missing set is controlled by an unusually rigid
three-parameter slice of the sextic coefficient space: the coefficients of
$w^5,w^4,w^3$ are fixed, so the first three power sums of the roots are fixed
as well. Write $Q_\star(w)=w^3-\frac12w^2-\frac18w-\frac5{16}$ for the cubic
of \eqref{f6:eq:uniquesquare}.

\begin{lemma}[At least three distinct roots \src{05}]\label{f6:kf:lem:threeroots}
Every polynomial $2w^6-2w^5-w^3+aw^2+bw+d$ has at least three distinct roots
over an algebraically closed field of characteristic zero.
\end{lemma}
\begin{proof}
For the monic polynomial $P/2$, Newton's identities give, with roots counted
with multiplicity, the power sums
\begin{equation}\label{f6:kf:eq:newton}
N_1=1,\qquad N_2=1,\qquad N_3=\frac52 .
\end{equation}
If all six roots were equal to $\alpha$, the first equation would give
$\alpha=1/6$, contradicting the second.

Suppose there are exactly two roots $\alpha,\beta$, of multiplicities
$m,6-m$. Interchange them so that $1\le m\le3$, and set
$\theta=\alpha-\beta$. The first moment gives
\[
\alpha=\frac16+\frac{6-m}{6}\theta,\qquad \beta=\frac16-\frac m6\theta .
\]
The second and third moments give, respectively,
\[
\theta^2=\frac5{m(6-m)},\qquad
\frac49+\frac{m(6-m)(3-m)}{18}\theta^3=\frac52 .
\]
Combining them yields
\begin{equation}\label{f6:kf:eq:tworootobstruction}
5(3-m)\theta=37 .
\end{equation}
For $m=3$ this is impossible. Squaring for $m=1$ or $m=2$ would require
$1369\,m(6-m)=125(3-m)^2$. The two sides are $6845$ and $500$ for $m=1$, and
$10952$ and $125$ for $m=2$. Neither equality holds.
\end{proof}

\begin{proposition}[The unique polynomial with no simple root \src{05}]
\label{f6:kf:prop:square}
The polynomial $2w^6-2w^5-w^3+aw^2+bw+d$ has no simple root if and only if
\begin{equation}\label{f6:kf:eq:star}
(a,b,d)=\Bigl(\frac{21}{32},\frac5{32},\frac{25}{128}\Bigr).
\end{equation}
At this parameter it equals $2Q_\star(w)^2$, and $Q_\star$ has three distinct
roots.
\end{proposition}
\begin{proof}[Second proof of Lemma~\ref{f6:lem:square} \src{05}]
If all roots have multiplicity at least two, Lemma~\ref{f6:kf:lem:threeroots}
forces exactly three distinct roots, each of multiplicity two. Write the
polynomial as $2(w^3+e_1w^2+e_2w+e_3)^2$. Comparison of the $w^5,w^4,w^3$
coefficients gives
\[
4e_1=-2,\qquad 2(e_1^2+2e_2)=0,\qquad 4(e_3+e_1e_2)=-1,
\]
so $e_1=-1/2$, $e_2=-1/8$ and $e_3=-5/16$. The three remaining coefficients
are exactly \eqref{f6:kf:eq:star}. Conversely, expansion proves the square
identity, and $\Disc Q_\star=-401/128\ne0$ shows that its three roots are
distinct.
\end{proof}

\Needspace{5\baselineskip}
\subsection{All seven multiplicity patterns \src{05}}
\label{f6:kf:sec:patterns}
Partitions of six with at least three parts are precisely the seven patterns
$1^6$, $2\,1^4$, $3\,1^3$, $2^2\,1^2$, $4\,1^2$, $3\,2\,1$, $2^3$ of
Theorem~\ref{f6:kf:thm:main}; Lemma~\ref{f6:kf:lem:threeroots} excludes the
others. A simple root contributes one point and every other root contributes
zero points (Theorem~\ref{f6:kf:thm:reconstruction}). This proves the counts
of Theorem~\ref{f6:kf:thm:main}(i) once all seven strata are shown to occur.

\begin{proposition}[Every listed pattern occurs \src{05}]\label{f6:kf:prop:witnesses}
Each of the seven patterns occurs in the family
$2w^6-2w^5-w^3+aw^2+bw+d$.
\end{proposition}
\begin{proof}
The following exact choices provide five patterns.
\begin{center}
\renewcommand{\arraystretch}{1.2}
\begin{tabular}{@{}lll@{}}
\toprule
$(a,b,d)$ & Polynomial or property & Pattern\\
\midrule
$(0,1,1)$ & discriminant $-1513772\ne0$ & $1^6$\\
$(1,0,0)$ & $w^2(w-1)(2w^3-1)$ & $2\,1^4$\\
$(0,0,0)$ & $w^3(2w^3-2w^2-1)$ & $3\,1^3$\\
$(-3/8,-1/8,-1/16)$ & $2(w^2+1/4)^2(w^2-w-1/2)$ & $2^2\,1^2$\\
$(21/32,5/32,25/128)$ & $2Q_\star(w)^2$ & $2^3$\\
\bottomrule
\end{tabular}
\end{center}
The displayed cofactors are squarefree and have no unwanted common roots, as
follows either by one derivative/gcd calculation or by direct substitution.

For the remaining patterns, prescribe a triple root $\rho$. Solving
$P(\rho)=P'(\rho)=P''(\rho)=0$ gives
\begin{align}
a(\rho)&=-30\rho^4+20\rho^3+3\rho,\notag\\
b(\rho)&=48\rho^5-30\rho^4-3\rho^2,\label{f6:kf:eq:triple}\\
d(\rho)&=-20\rho^6+12\rho^5+\rho^3.\notag
\end{align}
At these parameters
\begin{align}
P(w)&=(w-\rho)^3Q_\rho(w),\notag\\
Q_\rho(w)&=2w^3+(6\rho-2)w^2+(12\rho^2-6\rho)w+20\rho^3-12\rho^2-1,
\label{f6:kf:eq:triplefactor}
\end{align}
and in particular
\begin{equation}\label{f6:kf:eq:H4}
Q_\rho(\rho)=h_4(\rho):=40\rho^3-20\rho^2-1 .
\end{equation}
Choose a root of $h_4$. Neither $0$ nor $1/3$ is a root, and
$\partial_wQ_\rho(\rho)=10\rho(3\rho-1)\ne0$. The root $\rho$ therefore has
multiplicity exactly four in $P$. The other two roots cannot coincide, since
that would leave only two distinct roots. This realizes $4\,1^2$.

For the pattern $3\,2\,1$, compute
\begin{equation}\label{f6:kf:eq:K}
\Disc_wQ_\rho=-4\Lambda(\rho),\qquad
\Lambda=4320\rho^6-4320\rho^5+540\rho^4-448\rho^3+384\rho^2+36\rho+35 .
\end{equation}
This is a nonconstant polynomial and hence has a root over $k$. The
polynomials $\Lambda$ and $h_4$ are relatively prime. One way to see this
without a large Bézout identity is to use the fourfold-root case just
established: at a root of $h_4$, $Q_\rho$ has three distinct roots, one of
them $\rho$, because its remaining two roots are distinct and different from
$\rho$; thus $\Lambda$ does not vanish there. Choose a root of $\Lambda$.
Then $Q_\rho$ has a repeated root and does not vanish at $\rho$. It cannot be
a cube of one linear factor, since $P$ would have only two distinct roots.
Its pattern is therefore $2\,1$, giving the pattern $3\,2\,1$ for $P$.
\end{proof}

Every parameter triple in this proof corresponds to targets by
\eqref{f6:kf:eq:Yinv}. Algebraic witnesses for the one-point stratum suffice
for the assertion that the stratum is nonempty; no unproved assertion about
rational points on that stratum is needed. (The rows $(1,0,0)$,
$(-3/8,-1/8,-1/16)$ and $(21/32,5/32,25/128)$ are the parameters of the
four-, two- and zero-point targets of Table~\ref{f6:tab:witness}. The row
$(0,0,0)$ is a different three-point witness from that table's, whose
parameter is $(-7,15,-7)$. The row $(0,1,1)$ is the parameter of the second
arithmetic target $b_{**}$ of Section~\ref{f6:kf:sec:arithmetic}.)

\begin{theorem}[Exact omitted locus]\label{f6:thm:omitted}
Over an algebraically closed field, $\A^5\setminus\F(\A^5)$ is exactly
the surface $M$ of \eqref{f6:eq:M}. It is closed, smooth, irreducible, and
isomorphic to $\Gm\times\A^1$. An explicit parametrization, with
$c\ne0$ and $t\in k$, is
\begin{align}
 A&=-\frac{11}{32c},& B&=\frac{t}{c^2},\notag\\
 D&=\frac{t+25/128}{c^3},&
 E&=\frac{-t^2+(53/32)t+1773/4096}{c^4}.
 \label{f6:eq:Mparam}
\end{align}
\end{theorem}
\begin{proof}
There are no empty fibers on $c=0$ by \eqref{f6:eq:section}. For $c\ne0$,
Lemma~\ref{f6:lem:square} applies. The coefficients of $w^2,w,1$ in its
right-hand side are $21/32,5/32,25/128$. Comparing these with
\eqref{f6:eq:P} gives exactly \eqref{f6:eq:M}.

The first equation of \eqref{f6:eq:M} forces $c$ to be invertible even
scheme-theoretically: $c(-32A/11)=1$ in the coordinate ring. With
$t=c^2B$, the other equations solve uniquely to \eqref{f6:eq:Mparam}.
Thus the coordinate ring is $k[c,c^{-1},t]$. This proves all assertions,
including closedness in $\A^5$, dimension two, and codimension three.
\end{proof}

\begin{proof}[Second proof, through the effective parameters \src{05}]
No target with $c=0$ is omitted, by Section~\ref{f6:kf:sec:hyperplane}. For
$c\ne0$, Theorem~\ref{f6:kf:thm:reconstruction} and
Proposition~\ref{f6:kf:prop:square} show that a target is omitted exactly
when its effective parameters equal \eqref{f6:kf:eq:star}. Inserting these
constants into \eqref{f6:kf:eq:Yinv} gives
\[
Y_1=-\frac{11}{32},\quad Y_2=t,\quad Y_3=t+\frac{25}{128},\quad
Y_4=-t^2+\frac{53}{32}t+\frac{1773}{4096},
\]
and dividing by $c^i$ gives \eqref{f6:eq:Mparam}. The same set is cut out
globally by three polynomial equations, which are \eqref{f6:eq:M} cleared of
denominators:
\begin{align}
32cA+11&=0,\notag\\
128(c^3D-c^2B)-25&=0,\label{f6:kf:eq:Mideal}\\
32\bigl[c^4(E+B^2-AD)-2c^2B+cA\bigr]-5&=0.\notag
\end{align}
The first equation forces $c$ to be invertible, with $c^{-1}=-32A/11$. On
that open chart these equations are precisely $a=21/32$, $d=25/128$ and
$b=5/32$, and the coordinate change \eqref{f6:kf:eq:Yinv} proves that the
quotient ring is $k[c,c^{-1},t]$. This proves closedness in $\A^5$, not merely
closedness in $c\ne0$, and proves smoothness, the dimension, and that $M$ is
a complete intersection.
\end{proof}

\begin{remark}[Geometric image versus rational image]
The equality $\F(\A^5)=\A^5\setminus M$ is a statement about geometric
points. Over $\Q$, a surviving simple root need not be rational. There are
many rational targets outside $M$ with no rational preimage; the explicit
$S_6$ specialization in Section~\ref{f6:sec:galois} is one of them.
\end{remark}

\begin{remark}[An explicit rational omitted point \src{05}]
\label{f6:kf:cor:rationalmissing}
Source~05 states as a corollary that the rational target
\begin{equation}\label{f6:kf:eq:rationalmissing}
\Bigl(1,-\frac{11}{32},0,\frac{25}{128},\frac{1773}{4096}\Bigr)
\end{equation}
has no preimage even over $\C$; it is the zero-point witness of
Table~\ref{f6:tab:witness} ($c=1$, $t=0$ in \eqref{f6:eq:Mparam}). This is
an absolute geometric omission, not merely a failure to find rational
preimages. By contrast, Sections~\ref{f6:sec:galois}
and~\ref{f6:kf:sec:arithmetic} give rational targets with six complex
preimages but no rational preimage.
\end{remark}

\begin{corollary}[An obstruction to low-dimensional stabilization \src{05}]
\label{f6:kf:cor:nostabilization}
The map $\F$ is not polynomial-equivalent to
$\Phi\times\id_{\A^{5-m}}$ for any polynomial self-map $\Phi$ of $\A^m$
with $m\le3$, regardless of its geometric degree.
\end{corollary}
\begin{proof}
Polynomial equivalence means composition with polynomial automorphisms on the
source and on the target. A source automorphism does not change the image;
a target automorphism carries its complement isomorphically to the new
complement. For a product map that complement is
\[
\Omega_\Phi\times\A^{5-m},\qquad \Omega_\Phi=\A^m\setminus\im\Phi .
\]
If such an equivalence existed, this product would be isomorphic to the
nonempty irreducible surface $M$. The set $\Omega_\Phi$ is closed:
intersect the closed product complement with the section having identity
coordinates zero. Dimension then gives $\dim\Omega_\Phi=2-(5-m)=m-3$.
For $m<3$ this is impossible. For $m=3$, irreducibility forces the reduced
zero-dimensional variety $\Omega_\Phi$ to be one point, and the product
would be $\A^2$. But $M\simeq\Gm\times\A^1$ is not isomorphic to $\A^2$: its
coordinate ring has the nonconstant unit $c$, whereas every unit of
$k[y_1,y_2]$ is constant, and an isomorphism of varieties preserves the group
of units of its coordinate ring. This contradiction proves the claim and
Theorem~\ref{f6:kf:thm:main}(iv).
\end{proof}

The obstruction uses more than geometric degree: it excludes a stabilized
three-dimensional map even when that map also has geometric degree six.
\emph{A possible four-dimensional stabilization is not excluded by this
argument} (Section~\ref{f6:kf:sec:future}, direction~13). Source~05's
$N=\A^m\setminus\im\Phi$ is written $\Omega_\Phi$ here, since $\normal=\mathcal N$
is the normalization algebra and $\ord=\mathcal O$ the monogenic order.

\Needspace{5\baselineskip}
\subsection{Why codimension three is natural here \src{05}}
The family of fiber polynomials has three effective parameters $(a,b,d)$. At
the omitted parameter the sextic has three independent double roots.
Prescribing all three double-root degenerations uses all three transverse
parameter directions; the remaining target parameters $c$ and $t$ do not
occur in the fiber polynomial. That dimension count becomes a proof once
Proposition~\ref{f6:kf:prop:square} shows that no other all-multiple-root
pattern is possible.

The result does not say that codimension three is a universal maximum for
omitted loci of Keller maps; it says that this explicit map realizes it. Nor
should the codimension-three omitted set be confused with the nonproperness
set, which is a divisor (Section~\ref{f6:sec:discriminant}).

\Needspace{6\baselineskip}
\section{An explicit finite normalization}
\label{f6:sec:normalization}

\Needspace{5\baselineskip}
\subsection{The monogenic order is not normal}
Let
\[
 \base=k[c,A,B,D,E],\qquad
 \ord=\base[w]/(P_b(w)).
\]
Since $2$ is a unit, $\ord$ is free of rank six over $\base$.
The polynomial $P_b$ is irreducible: as a polynomial in $E$, its leading
coefficient is $c^4w$, and its constant term is divisible by neither $c$
nor $w$. Gauss's lemma gives the assertion, as in the proof of
Theorem~\ref{f6:thm:chart}. Thus $\ord$ is a domain.

The source function field is its fraction field. Indeed, away from $c=0$
the root chart identifies $w=cv$ and gives the full source coordinates
rationally in $b,w$. Nevertheless, $\ord$ is not normal along its
collapsed locus $c=w=0$. At $c=0$ its defining equation is
\[
 P_b(w)=w^2g(w),
\]
independent of $A,B,D,E$. That single double root has forgotten the
quadratic $v^2+Av-B$ which controls the two additional source branches.

\Needspace{5\baselineskip}
\subsection{One fraction and six basis elements}
Inside $\operatorname{Frac}(\ord)$ put
\begin{equation}\label{f6:eq:zeta}
 \zeta=\frac{wJ(w)}{c},\qquad \normal=\base[w,\zeta].
\end{equation}
The polynomial relation $P_b=0$ gives the symmetric identities
\begin{equation}\label{f6:eq:rankone}
 wJ=c\zeta,\qquad w(\zeta-T)=cL,\qquad
 \zeta(\zeta-T)=LJ.
\end{equation}
They are exactly the multiplication rules \eqref{f6:eq:rules}.

\begin{theorem}[Finite free normalization]\label{f6:thm:normalization}
The algebra $\normal$ is the integral closure of both $\ord$ and $\base$
in the source function field. It is free over $\base$ with basis
\[
 \mathcal B=(1,w,w^2,w^3,w^4,\zeta).
\]
The affine variety $Z=\Spec\normal$ is smooth over $k$, and
$\pi:Z\to\A^5$ is finite flat of degree six.
\end{theorem}
\begin{proof}
\emph{Finite freeness.} The $\base$-span of $\mathcal B$ contains $1,w$
and $\zeta$. It is stable under multiplication by $w$ and by $\zeta$:
use the three rules \eqref{f6:eq:rules}, reducing $w^5$, $w\zeta$, and
$\zeta^2$. This span is therefore the whole algebra $\normal$.
The six elements are independent, because relative to
$1,w,\ldots,w^5$ the last element has coefficient $2/c$ on $w^5$.
The degree-six irreducibility of $P_b$ gives independence over
$\operatorname{Frac}(\base)$, and hence over $\base$.
Thus $\normal$ is finite free of rank six. The same reduction and
independence show that the three displayed relations form a presentation:
any polynomial in $w,\zeta$ reduces to the basis span, and no additional
relation can remain there. In particular its elements are integral over
$\base$.

\emph{Three useful charts.} On $c\ne0$ we have
$\normal[1/c]=\ord[1/c]$. On $w\ne0$, the second relation in
\eqref{f6:eq:rankone} gives $\zeta=T+cL/w$, so
$\normal[1/w]=\ord[1/w]$. On $J\ne0$ put $v=\zeta/J$. The first
relation gives $w=cv$, and the third gives
\[
 v^2J(cv)-Tv-L=0.
\]
Conversely $w=cv$ and $\zeta=vJ(cv)$ satisfy all three relations. Therefore
\begin{equation}\label{f6:eq:Jchart}
 \normal[1/J]\simeq
 \base[v,1/J(cv)]/(q_b(v)).
\end{equation}
The three opens cover $Z$: a point with $c=w=0$ has $J=1$.

\emph{Smoothness on those charts.} On $c\ne0$, work in the hypersurface
$P_b=0$. Its derivative with respect to $E$ is $c^4w$. If that derivative
vanishes, then $w=0$, and the derivative with respect to $D$ is $c^3$,
which is nonzero. Thus this open is smooth. At a point with $c=0,w\ne0$,
the equation forces $g(w)=0$ and
\[
 \frac{\partial P_b}{\partial w}=w^2g'(w)\ne0,
\]
since $g=(w-1)(2w^3-1)$ has four distinct nonzero roots.
Finally, at a point with $c=w=0$, chart \eqref{f6:eq:Jchart} applies and
$\partial q_b/\partial B=-1$. These checks prove smoothness everywhere.

A smooth integral algebra over a characteristic-zero field is normal.
The finite integral algebra $\normal$ has the same fraction field as
$\ord$. Any element of that field integral over $\ord$ is integral over
$\normal$ and therefore belongs to $\normal$. This proves that
$\normal$ is its integral closure. Since $\ord$ is itself integral over
$\base$, it is also the integral closure of $\base$ in that field.
\end{proof}

\Needspace{5\baselineskip}
\subsection{The exact fiber of the finite model at \texorpdfstring{$c=0$}{c=0}}
Reducing \eqref{f6:eq:rankone} modulo $c$ gives
\[
 wg(w)=0,\qquad w(\zeta+A)=0,\qquad
 \zeta(\zeta+A)=Bg(w).
\]
As $g(0)=1$, the zero locus of $w$ and the zero locus of $g$ are disjoint.
The Chinese remainder decomposition is
\begin{equation}\label{f6:eq:CRT}
 \normal/(c)\simeq
 k[A,B,D,E,\zeta]/(\zeta^2+A\zeta-B)
 \ \times\ k[A,B,D,E,w]/(g(w)).
\end{equation}
On the second factor $\zeta=-A$; on the first $w=0$.
The fiber has length two from the quadratic factor and length four from
the residual quartic factor. This accounts for all six sheets, even where
the quadratic becomes nonreduced.

\begin{theorem}[The source as an open part of its finite model]
\label{f6:thm:boundary}
The map $\F$ factors as an open immersion followed by the finite map:
\[
 \A^5\ \xhookrightarrow{\ j\ } Z\ \xrightarrow{\ \pi\ }\A^5.
\]
Over an algebraically closed field, let $\rho_1,\rho_2,\rho_3$ be the
three roots of $2w^3=1$, and let
\[
 E_i=\{c=0,\ w=\rho_i,\ \zeta=-A\}\subset Z.
\]
Each $E_i$ is isomorphic to $\A^4$, and $\pi$ maps it isomorphically onto
$c=0$. The exact complement of the source is
\begin{equation}\label{f6:eq:boundary}
 Z\setminus j(\A^5)=\Ram(\pi)\ \cup\ E_1\cup E_2\cup E_3.
\end{equation}
The three $E_i$ are unramified and disjoint from $\Ram(\pi)$.
\end{theorem}
\begin{proof}
On the source define $w=\gamma u_1$. The rational function $\zeta$ of
\eqref{f6:eq:zeta} is regular there. On $x\ne0$ it equals $vJ(cv)$ by the
root chart. On $w\ne0$ it equals $T+cL/w$. These opens cover the source,
because $x=0$ implies $\gamma=u_1=w=1$. Thus they give a morphism
$j:\A^5\to Z$ over the target.

For $c\ne0$, $\pi$ is the sextic root cover. Its unramified points are
exactly the simple roots, and Theorem~\ref{f6:thm:chart} identifies all of
these with source points. For $c=0,w=0$, decomposition
\eqref{f6:eq:CRT} and chart \eqref{f6:eq:Jchart} identify the unramified points
with the simple roots of $v^2+Av-B$; all of these are again source points.
The root $w=1$ at $c=0$ is unramified, and is precisely the image of the
source hyperplane under \eqref{f6:eq:section}. The other three quartic roots
are also unramified, but cannot come from the source: $x=0$ gives $w=1$,
whereas $x\ne0$ and $c=0$ give $\gamma=0$ and $w=0$.

This proves the asserted set-theoretic image and shows that $j$ is
injective on geometric points. It lands in the \'etale locus of $\pi$.
Both its source and that locus are \'etale over the target; a morphism
between such schemes is \'etale \cite[Tag 02GH]{StacksEtale}. Geometric
injectivity makes it universally injective. An \'etale universally
injective morphism is an open immersion \cite[Tag 025G]{StacksOpen}.
Finally the four nonzero roots of $g$ are simple everywhere over $c=0$,
so all three removed components are unramified and disjoint from the
ramification locus. The description of each $E_i$ follows directly from
\eqref{f6:eq:CRT}.
\end{proof}

The phrase ``three escapes'' refers precisely to these three missing,
unramified sheets over the generic point of $c=0$. At a generic point of
the other exceptional component, the mechanism is different: two sheets
meet at one ramification point, and that point is removed from the source.

\Needspace{6\baselineskip}
\section{The discriminant and the nonproperness locus}
\label{f6:sec:discriminant}

\Needspace{5\baselineskip}
\subsection{Removing a spurious discriminant factor}
Let
\begin{equation}\label{f6:eq:universaldisc}
 \Delta(a,b,d)=\Disc_w(2w^6-2w^5-w^3+aw^2+bw+d).
\end{equation}
It is unnecessary to print its full expansion. The target substitution is
\begin{equation}\label{f6:eq:discsub}
 a=1+cA,\qquad b=-cT=c^4\kk-2c^2B+cA,
 \qquad d=-c^2L=c^3D-c^2B.
\end{equation}

\begin{theorem}[Normalized discriminant]\label{f6:thm:H}
The expression
\begin{equation}\label{f6:eq:H}
 H=\frac{\Delta(1+cA,-cT,-c^2L)}{c^2}
\end{equation}
is a polynomial, irreducible over every algebraically closed field of
characteristic zero, and
\begin{equation}\label{f6:eq:H0}
 H|_{c=0}=-108(A^2+4B).
\end{equation}
In the basis $\mathcal B$ of Theorem~\ref{f6:thm:normalization}, the trace
pairing of $\normal/\base$ has determinant $H/256$.
\end{theorem}
\begin{proof}
The discriminant of the power basis of the monic polynomial $P_b/2$ is
$\Disc(P_b)/2^{10}$. The change from that power basis to $\mathcal B$
has determinant $2/c$. A trace discriminant changes by the square of
the basis determinant. Therefore
\begin{equation}\label{f6:eq:discbasis}
 \Disc(\mathcal B)=\frac{4}{c^2}\frac{\Disc(P_b)}{2^{10}}
                 =\frac{H}{256}.
\end{equation}
The left-hand side belongs to $\base$ because all multiplication rules
have polynomial coefficients. This proves polynomiality of $H$.

For its value at $c=0$, first assume $A^2+4B\ne0$. The two roots of $P_b$
near zero have the form $cv_++O(c^2)$ and $cv_-+O(c^2)$, where
$v_\pm$ are the two roots of $v^2+Av-B$. Their squared difference
contributes $c^2(A^2+4B)$ to first order. The other four roots tend to
the distinct roots of $g$. The leading-coefficient factors and all
cross-differences combine to $\Disc(g)=-108$: explicitly, the product
of the four roots of $g$ is $1/2$, the cross-difference contribution is
$(1/2)^4$, and the leading-coefficient ratio is $2^{10}/2^6=16$.
Thus the coefficient of $c^2$ in $\Disc(P_b)$ is
$-108(A^2+4B)$. This formal expansion is valid on a dense open set and
therefore gives the polynomial identity \eqref{f6:eq:H0}. It also follows
by direct substitution into the exact resultant.

For irreducibility, the double-root incidence of the polynomial in
\eqref{f6:eq:universaldisc} is parametrized by $(r,a)$:
\begin{align}
 b&=-12r^5+10r^4+3r^2-2ar,\notag\\
 d&=10r^6-8r^5-2r^3+ar^2.\label{f6:eq:incidence}
\end{align}
Its image is exactly the discriminant hypersurface; hence that
hypersurface is irreducible. The discriminant is reduced, not a higher
power: at $(a,b,d)=(1,0,0)$ the polynomial is
$w^2(w-1)(2w^3-1)$ with one ordinary double root, and varying $d$ gives
a simple zero of the discriminant. For example the coefficient of $d$
in $\Delta(1,0,d)$ is $432\ne0$.

After inverting $c$, substitution \eqref{f6:eq:discsub} is a coordinate
isomorphism
\[
 (c,A,B,D,\kk)\longleftrightarrow(c,B,a,b,d).
\]
Therefore $H$ is irreducible after inverting $c$. Any additional factor
before localization would have to be a power of $c$. Equation
\eqref{f6:eq:H0} excludes such a factor. This proves irreducibility.
\end{proof}

For a finite locally free algebra in characteristic zero, the trace
pairing is nondegenerate precisely on the \'etale locus
\cite[Tag 0BJF]{StacksDisc}. Thus $V(H)$ is the branch locus of the
finite normalization, while the factor $c^2$ in the unnormalized sextic
discriminant records its nonnormal collapse. The supplied sparse file
stores $H$ with $179$ monomials in $(c,A,B,D,\kk)$.

\Needspace{5\baselineskip}
\subsection{The discriminant surface and its normalization \src{05}}
Let $\dd=V(\Delta)\subset\A^3_{a,b,d}$ be the discriminant surface in
coefficient space. Writing it as a discriminant fixes its scalar
normalization; no choice of a reduced defining equation is hidden in later
constants. Source~05 ships the expanded $54$-term polynomial $\Delta$
(Appendix~\ref{f6:app:reproduce}). The incidence \eqref{f6:eq:incidence}
used for irreducibility in Theorem~\ref{f6:thm:H} is in fact the
normalization of $\dd$, with $r$ written $\rho$:

\begin{proposition}[An explicit normalization of $\dd$ \src{05}]
\label{f6:kf:prop:Dnorm}
The polynomial $\Delta$ is absolutely irreducible up to a nonzero scalar. The
normalization of $\dd$ is $\A^2_{a,\rho}$, with finite normalization map
\begin{equation}\label{f6:kf:eq:Dnorm}
(a,\rho)\longmapsto
\bigl(a,\ -12\rho^5+10\rho^4+3\rho^2-2a\rho,\
       10\rho^6-8\rho^5-2\rho^3+a\rho^2\bigr).
\end{equation}
The points above a parameter are exactly its distinct repeated roots.
\end{proposition}
\begin{proof}
Solving $P(\rho)=P'(\rho)=0$ for $b,d$ gives \eqref{f6:kf:eq:Dnorm}. Thus the
repeated-root incidence variety is the affine plane with coordinates
$(a,\rho)$. Its projection is finite: $\rho$ is integral over the
coefficient ring, since it satisfies the monic equation $P(\rho)/2=0$. The
image is exactly the discriminant zero set, so that set is irreducible.

At $(a,b,d)=(1,0,0)$ the polynomial has exactly one double root and four
simple roots, and direct differentiation of the discriminant gives
$\partial_d\Delta(1,0,0)=432\ne0$. Thus $\Delta$ is reduced along its unique
irreducible component: a polynomial with irreducible zero set is a scalar
times a power of the irreducible defining polynomial, and the nonzero
derivative forces that power to be one.

The condition $\deg\gcd(P,P')=1$ is open in the discriminant, by the rank
conditions on the Sylvester matrix. It holds at the displayed witness and is
therefore dense. On this locus the projection has one-point fibers, so in
characteristic zero it is birational. A finite birational map from the
normal affine plane is the normalization. The description of its fibers
follows from its incidence definition.
\end{proof}

For example, the normalization has three distinct points above the omitted
parameter \eqref{f6:kf:eq:star}. A single irreducible discriminant surface can
therefore have three local analytic branches at that parameter: global
irreducibility and local reducibility are different notions.

\begin{proof}[Second proof of \eqref{f6:eq:H0} and of the irreducibility of
$H$, by coprime factor lifting \src{05}]
At $c=0$ the sextic is $w^2g(w)$, with $g(0)=1$ and $\Disc g=-108$. Work in
the formal series ring $k[A,B,D,E][[c]]$. Coprime factor lifting splits $P_b$
uniquely into a monic quadratic factor $\phi_c$ reducing to $w^2$ and a
degree-four factor $\psi_c$ reducing to $g$; only their initial terms are
needed. The effective parameters have expansions
\[
a=1+cA,\qquad b=cA-2c^2B+O(c^4),\qquad d=-c^2B+c^3D .
\]
The quadratic factor has linear coefficient $cA+O(c^2)$ and constant
coefficient $-c^2B+O(c^3)$, so $\Disc\phi_c=c^2(A^2+4B)+O(c^3)$. The product
formula for discriminants gives
$\Disc P_b=(\Disc\phi_c)(\Disc\psi_c)\Res(\phi_c,\psi_c)^2$, and at $c=0$ the last two
factors are $-108$ and $1$. Consequently
\begin{equation}\label{f6:kf:eq:Dleading}
\Delta(a,b,d)=-108c^2(A^2+4B)+O(c^3).
\end{equation}
As an identity about the first three coefficients of a polynomial in $c$,
this proves divisibility by $c^2$ and shows that $c^3$ does not divide it;
the supplied script checks the same coefficient identity directly. After
inverting $c$, the coefficient change identifies the ring with
$k[c,c^{-1},t,a,b,d]$, and Proposition~\ref{f6:kf:prop:Dnorm} shows that the
substituted discriminant is irreducible there. Any factorization of $H$ in
the polynomial ring would therefore have a factor that becomes a unit after
inverting $c$, hence a scalar times a power of $c$, which
\eqref{f6:kf:eq:Dleading} rules out.
\end{proof}

\Needspace{5\baselineskip}
\subsection{Two components of nonproperness}
\begin{definition}
For a polynomial map $f:\C^n\to\C^n$, its nonproperness set consists
of targets $b$ for which no neighborhood $V$ of $b$ has
$f^{-1}(V)\to V$ proper. Equivalently, there exists an unbounded sequence
of source points whose images converge to $b$.
\end{definition}

\begin{theorem}[Exact nonproperness locus]\label{f6:thm:nonproper}
For $\F:\C^5\to\C^5$,
\begin{equation}\label{f6:eq:nonproper}
 \NP(\F)=V(cH)=V(c)\cup V(H).
\end{equation}
These are two distinct irreducible hypersurfaces. Their intersection is
\begin{equation}\label{f6:kf:eq:boundaryintersection}
 c=0,\qquad A^2+4B=0.
\end{equation}
The generic fiber on $V(c)$ has three points; the generic fiber on
$V(H)$ has four points; every point of their intersection has one point
in its original fiber. \textup{(Source~05:)} Both components are smooth
along \eqref{f6:kf:eq:boundaryintersection} and meet transversely there;
smoothness of $V(H)$ is asserted along that intersection, not everywhere.
\end{theorem}
\begin{proof}
If $cH\ne0$, the sextic has six distinct roots and the original map has
six points over the target. Since $\F$ has nonzero Jacobian, choose six
disjoint local inverse neighborhoods. On a sufficiently small common
target neighborhood they supply six inverse branches. There can be no
additional branch, because every fiber has at most six points. After
shrinking to a closed ball inside this neighborhood, the preimage is
a finite union of compact continuous images of that ball. Thus $\F$ is
proper near the target.

At a target on $V(cH)$, the fiber has fewer than six points by
Corollary~\ref{f6:cor:count}. Suppose the map were proper over a small
connected neighborhood. A proper local homeomorphism with finite fibers
is a covering onto a union of components; its fiber cardinality is
locally constant. Nearby points with $cH\ne0$ have six preimages, a
contradiction. This proves \eqref{f6:eq:nonproper}.

Irreducibility and the intersection statement follow from
Theorem~\ref{f6:thm:H}. At a generic point of the discriminant there is one
double root and four simple roots, so the original fiber has four
points. The hyperplane and intersection counts are
\eqref{f6:eq:c0count}. For transversality \src{05}: by \eqref{f6:eq:H0},
$\partial H/\partial B=-432$ on $c=0$, which is nonzero at every point of
\eqref{f6:kf:eq:boundaryintersection}. The gradients of $c$ and $H$ are
therefore linearly independent there, proving smoothness and
transversality.
\end{proof}

\begin{proof}[Second proof of \eqref{f6:eq:nonproper}, through the finite
completion \src{05}]
On $U=\{c\ne0,\ \Delta\ne0\}$ the sextic is squarefree. In the rank-six
algebra of \eqref{f6:kf:eq:completion}, $P'$ is invertible, because a Bézout
identity between $P$ and $P'$ exists after inverting their resultant. The
entire finite completion is therefore the actual source over $U$, so
$\F^{-1}(U)\to U$ is finite \'etale of degree six, and proper. At every point
outside $U$ the fiber-size classification gives fewer than six points: one
or three on $c=0$, and at most four on $c\ne0$, $\Delta=0$. Suppose the map
were proper over a neighborhood of such a point. An everywhere local
biholomorphism that is proper over a sufficiently small connected
neighborhood is a finite covering there; its number of sheets equals the
size of its central fiber and is locally constant. The neighborhood also
contains points of $U$, whose fibers have six points, a contradiction. This
argument includes omitted points: a proper covering with zero sheets would
have empty image on the neighborhood.
\end{proof}

The four nearby fiber counts at a point of the intersection
\eqref{f6:kf:eq:boundaryintersection} are worth distinguishing \src{05}:
\[
\begin{array}{c|c}
\text{location}&\text{fiber size}\\\hline
c\ne0,\ H\ne0&6\\
c=0,\ H\ne0&3\\
c\ne0,\ H=0\text{ near the intersection}&4\\
c=0,\ H=0&1
\end{array}
\]
The third row holds because $V(H)$ is smooth near the intersection, so the
off-hyperplane sextic has exactly one double root there. In this local
picture three sheets are lost on one divisor and two on the other; the one
visible hyperplane point remains at their intersection.

There is an equivalent finite-model explanation. The image under $\pi$
of the ramification locus is $V(H)$, whereas the three deleted
unramified components all map to $V(c)$. Equation \eqref{f6:eq:boundary}
therefore describes exactly which points of the finite cover have to be
added to restore properness. The omitted surface $M$ is much smaller
than the nonproperness set: points of $M$ lose every sheet, while a
generic nonproper target retains several sheets.

\Needspace{6\baselineskip}
\subsection{Every local analytic type of the discriminant \src{05}}
\label{f6:kf:sec:local}
The three-parameter slice $(a,b,d)$ is not accidentally too small to see the
relevant degenerations: it is transverse to every contact pattern that
actually occurs. This gives a local analytic description at every point of
$\dd$, not only at the omitted parameter. (This partly answers source~02's
question~3 in Section~\ref{f6:sec:future}.)

\begin{lemma}[Hermite transversality \src{05}]\label{f6:kf:lem:hermite}
Suppose $P=2w^6-2w^5-w^3+a_0w^2+b_0w+d_0$ has distinct repeated roots
$\rho_1,\dots,\rho_j$ of multiplicities $m_1,\dots,m_j$, and put
$\ell=\sum_{i=1}^j(m_i-1)$. Then $\ell\le3$, and the perturbations
$\delta a\,w^2+\delta b\,w+\delta d$ map surjectively to
\begin{equation}\label{f6:kf:eq:hermite}
\bigoplus_{i=1}^j\C[w]/(w-\rho_i)^{m_i-1}.
\end{equation}
\end{lemma}
\begin{proof}
The sum of multiplicity excesses is six minus the number of distinct roots,
so Lemma~\ref{f6:kf:lem:threeroots} gives $\ell\le3$. The Chinese remainder
theorem identifies \eqref{f6:kf:eq:hermite} with
$\C[w]/\prod_i(w-\rho_i)^{m_i-1}$. The modulus has degree $\ell\le3$, so
every residue class has a representative of degree at most two. This is
precisely the asserted surjectivity.
\end{proof}

\begin{theorem}[Local analytic normal forms \src{05}]\label{f6:kf:thm:localtypes}
In suitable local analytic coordinates $(\sigma_1,\sigma_2,\sigma_3)$ on the
coefficient space $\A^3_{a,b,d}$, $\Delta$ is a nonvanishing analytic unit
times the equation in the following table; unused coordinates are smooth
parameters.
\begin{center}
\renewcommand{\arraystretch}{1.25}
\begin{tabular}{@{}lll@{}}
\toprule
Multiplicity pattern & Local equation for $\dd$ & Fiber size\\
\midrule
$1^6$ & no discriminant point & $6$\\
$2\,1^4$ & $\sigma_1$ & $4$\\
$3\,1^3$ & $4\sigma_1^3+27\sigma_2^2$ & $3$\\
$2^2\,1^2$ & $\sigma_1\sigma_2$ & $2$\\
$4\,1^2$ & $\Delta_4(\sigma_1,\sigma_2,\sigma_3)$ & $2$\\
$3\,2\,1$ & $\sigma_3(4\sigma_1^3+27\sigma_2^2)$ & $1$\\
$2^3$ & $\sigma_1\sigma_2\sigma_3$ & $0$\\
\bottomrule
\end{tabular}
\end{center}
Here $\Delta_4$ is the discriminant of $\omega^4+\sigma_1\omega^2+\sigma_2\omega+\sigma_3$,
\begin{equation}\label{f6:kf:eq:quarticdisc}
\Delta_4=256\sigma_3^3-128\sigma_1^2\sigma_3^2+144\sigma_1\sigma_2^2\sigma_3
 -27\sigma_2^4+16\sigma_1^4\sigma_3-4\sigma_1^3\sigma_2^2 .
\end{equation}
In particular, the singular points of $\dd$ are exactly the parameters with a
root of multiplicity at least three or with at least two distinct repeated
roots.
\end{theorem}
\begin{proof}
Choose disjoint small disks about the distinct roots of the reference
polynomial. In a neighborhood of its coefficients, analytic factorization
groups the roots in each disk into a monic factor of the corresponding
degree. This follows, for example, from the implicit function theorem for
multiplication of coprime monic factors, whose linearization is invertible
by the polynomial Chinese remainder theorem. The resultants between
different factors do not vanish in this neighborhood, so the discriminant
product formula expresses $\Delta$, up to an analytic unit, as the product of
the discriminants of the repeated-root factors; simple-root factors
contribute only units.

Center the degree-$m_i$ factor about its mean root to remove its
degree-$(m_i-1)$ term. Its remaining deformation coefficients have tangent
space $\C[\omega]/\omega^{m_i-1}$. The differential from the original
perturbation $h(w)$ to these coefficients is the class of $h$ at $\rho_i$,
multiplied by the inverse of the other factors; this multiplier is a unit in
the corresponding truncated local ring. The combined differential is
therefore surjective by Lemma~\ref{f6:kf:lem:hermite}. Complete these $\ell$
deformation coefficients to three analytic coordinates by the inverse
function theorem.

Thus the repeated-root factors, after centering, are independent universal
monic polynomials of degrees $m_i$ with their degree-$(m_i-1)$ coefficients
removed. A quadratic gives one coordinate as its discriminant; a centered
cubic gives $-4\sigma_1^3-27\sigma_2^2$; a centered quartic gives
\eqref{f6:kf:eq:quarticdisc}. Multiplying these factors yields the table,
with harmless nonzero constants absorbed into units. In the smooth
double-root case the differential is nonzero; all other displayed equations
have zero differential at their base point. This proves the singular-locus
criterion as well.
\end{proof}

\Needspace{5\baselineskip}
\subsection{Three transverse escape branches along the omitted surface
\src{05}}
\begin{corollary}[Normal crossings and local monodromy \src{05}]
\label{f6:kf:cor:crossings}
Near any point of $M$ there are analytic coordinates
$(\epsilon_1,\epsilon_2,\epsilon_3,\lambda_1,\lambda_2)$ on the target for
which
\[
V(H)=V(\epsilon_1\epsilon_2\epsilon_3),\qquad M=V(\epsilon_1,\epsilon_2,\epsilon_3).
\]
If exactly $j$ of the $\epsilon_i$ vanish, the fiber has $6-2j$ points. A
small loop around $\epsilon_i=0$ exchanges one pair of sheets and fixes the
other four. The local monodromy image is $(\Z/2\Z)^3$, acting on three
disjoint pairs.
\end{corollary}
\begin{proof}
At the omitted parameter, $P=2Q_\star^2$ and $Q_\star$ has three distinct
roots. The $2^3$ row of Theorem~\ref{f6:kf:thm:localtypes} gives the product
equation; the two extra coordinates come from $c$ and $t$
(Section~\ref{f6:kf:sec:parameters}, $c\ne0$ on $M$). More explicitly, each
of the three quadratic root clusters can be centered and written as
$\omega_i^2-\epsilon_i$. Its two roots are distinct and survive
reconstruction when $\epsilon_i\ne0$; when $\epsilon_i=0$, its double root is
entirely removed by localization at $P'$. The three clusters stay disjoint,
so exactly two points are lost for each zero coordinate. Analytic
continuation of $\sqrt{\epsilon_i}$ around zero exchanges the two roots in
that cluster and no others. These three transpositions commute and are
independent.
\end{proof}

This also explains why the empty fiber is not produced by six source points
colliding at one affine source point. The differential of $\F$ is invertible
everywhere. The collisions occur in the finite completion
\eqref{f6:kf:eq:completion} (equivalently in $Z$), and the collision points
have been removed from the affine source: with $c$ tending to a nonzero
value, already the coordinate $x=c/P'(w)$ diverges as $P'(w)$ tends to zero.
(This matches source~02's question~4, which asks for the escape rates.)

Together, Theorems~\ref{f6:thm:nonproper} and~\ref{f6:kf:thm:localtypes} give
all local analytic types of the nonproperness divisor. Off $c=0$ they are
the types in the table, with two smooth factors. On $c=0$ the divisor is
smooth away from \eqref{f6:kf:eq:boundaryintersection}, and has two
transverse branches along that intersection.

\Needspace{6\baselineskip}
\section{A general quadratic-contraction normalization lemma}
\label{f6:sec:general}
The short normalization proof is not peculiar to the large coefficients
in the stage. Its algebraic core applies to a wider class of degenerating
root covers. The following lemma records that mechanism without claiming
that every such cover is a Keller map.

\begin{theorem}[Quadratic-contraction algebra]\label{f6:thm:general}
Let $A_0$ be a domain, let $c\in A_0$ be nonzero, and let
$J(w)\in A_0[w]$ have degree $m\ge1$ and leading coefficient
$a_m\in A_0^\times$. Choose $T,L\in A_0$ and suppose
\[
 P(w)=w^2J(w)-cTw-c^2L
\]
is irreducible over $\operatorname{Frac}(A_0)$. In its root field define
$\zeta=wJ(w)/c$ and $N=A_0[w,\zeta]$. Then:
\begin{enumerate}[label=\textup{(\roman*)}]
\item $N$ is finite free of rank $m+2$ over $A_0$, with basis
$1,w,\ldots,w^m,\zeta$ and relations
\[
 wJ=c\zeta,\qquad w\zeta=Tw+cL,\qquad
 \zeta^2=T\zeta+LJ.
\]
\item Its trace discriminant in that basis equals
\[
 \frac{\Disc(P)}{a_m^{2m}c^2}.
\]
\item On $c\ne0$ or $w\ne0$ the algebra equals the original monogenic
order. On $J\ne0$ it has the chart
\[
 A_0[v,1/J(cv)]/(v^2J(cv)-Tv-L).
\]
If $c$ and $J(0)$ generate the unit ideal, these three opens cover
$\Spec N$. If their rings are normal, $N$ is the integral closure of
the monogenic order in its fraction field.
\end{enumerate}
\end{theorem}
\begin{proof}
The three relations follow by multiplying $P=0$ by suitable powers of
$J/c$ and cancelling in the fraction field. They show that the span of
the displayed basis is closed under multiplication by $w$ and $\zeta$;
the leading coefficient $a_m$ is a unit, so $w^{m+1}$ reduces using
$wJ=c\zeta$. Independence follows from the coefficient $a_m/c$ of
$w^{m+1}$ in $\zeta$. This proves (i).

The monic power-basis discriminant is
$\Disc(P)/a_m^{2m+2}$. The basis determinant is $a_m/c$, giving (ii).
The chart calculations are the same elementary substitutions as in
\eqref{f6:eq:Jchart}. If $c=w=J=0$, then $c=J(0)=0$, which is excluded by
the unit-ideal assumption. Thus the charts cover. Normality is local;
a finite normal algebra in the same fraction field is the integral
closure. This proves (iii).
\end{proof}

For $F_6$, take $m=4$, $a_m=2$, and $J(0)=1+cA$.
The smoothness checks in Theorem~\ref{f6:thm:normalization} supply the
normality hypotheses. For a general $J,T,L$, those hypotheses must be
checked: the finite free algebra in this lemma is not automatically
normal. This distinction is important when attempting degree-twelve or
higher analogues.

The lemma explains the otherwise mysterious square factor $c^2$ in a
root discriminant: one basis vector must be divided by $c$ to separate
the two collapsed directions. It also provides a practical alternative
to normalizing a large graph ideal in all source and target variables.

\Needspace{6\baselineskip}
\section{Geometric monodromy and an arithmetic radical obstruction}
\label{f6:sec:galois}

\Needspace{5\baselineskip}
\subsection{A Morse slice with full symmetric monodromy}
Set
\[
 p(w)=2w^6-2w^5-w^3+w^2.
\]
On the target line $(c,A,B,D,E)=(1,0,0,t,0)$, the fiber polynomial is
$p(w)+t$. Its discriminant is
\begin{equation}\label{f6:eq:morsedisc}
 \Disc_w(p(w)+t)
 =-4t(32t+3)(11664t^3-5680t^2+1149t-36).
\end{equation}
The polynomial on the right is squarefree. This can be checked by the
Euclidean algorithm, and is independently verified in the code.

\begin{theorem}[Geometric $S_6$]\label{f6:thm:geometricS6}
The six-sheeted cover over $\C^5\setminus V(cH)$ has monodromy group
$S_6$. Equivalently, $P_b$ has Galois group $S_6$ over
$\C(c,A,B,D,E)$.
\end{theorem}
\begin{proof}
For a polynomial $p$ of degree six, the discriminant of $p(w)+t$ is a
nonzero constant times the product of $t+p(r)$ over the five critical
points $r$ of $p$, counted with multiplicity. The squarefreeness of
\eqref{f6:eq:morsedisc} implies that all five critical points are simple
and their critical values are distinct.

Remove the five critical values from the target $t$-line. The inverse
image under $t=-p(w)$ is the complex $w$-line with finitely many points
removed, hence is connected. Its degree-six covering therefore has
transitive monodromy. A small loop around a simple critical value
interchanges exactly two sheets, because the local map has a nonzero
quadratic term. The fundamental group of the punctured line is generated
by these small loops. Thus its monodromy is a transitive permutation
group generated by transpositions.

Place an edge between two sheet labels whenever a generating
transposition exchanges them. Transitivity makes this graph connected.
Transpositions along a connected graph generate all permutations, so
the group is $S_6$. Restricting the full target cover to this line cannot
enlarge its monodromy. Since the restricted group is already $S_6$, the
full group is exactly $S_6$.
\end{proof}

\Needspace{5\baselineskip}
\subsection{A second Morse slice \src{05}}
Source~05 proves the same theorem on a different line. Local crossings alone
(Corollary~\ref{f6:kf:cor:crossings}) do not determine the global permutation
group; in this example a one-parameter slice supplies all permutations. Set
$a=0$, $b=1$ and let $d=\tau$ vary. On the target line
\begin{equation}\label{f6:kf:eq:line}
(c,A,B,D,E)=(1,-1,0,\tau,2-\tau)
\end{equation}
(that is, $c=1$, $t=0$ in \eqref{f6:kf:eq:Yinv}) the fiber polynomial is
\begin{equation}\label{f6:kf:eq:P}
\hat p(w)+\tau,\qquad \hat p(w)=2w^6-2w^5-w^3+w .
\end{equation}
Two exact identities are
\begin{align}
\hat p'(w)&=(w-1)(2w-1)(6w^3+4w^2+3w+1),\label{f6:kf:eq:Pprime}\\
\Disc_w(\hat p(w)+\tau)&=-4\tau(32\tau+11)
 (11664\tau^3-6976\tau^2+5485\tau-1372).\label{f6:kf:eq:criticalvalues}
\end{align}
The last polynomial is squarefree. For instance, its cubic factor has
discriminant
\[
-3990738242362176\ne0,
\]
and its values at $0$ and $-11/32$ are $-1372$ and $-9329787/2048$. Hence $\hat p$ has five simple critical points
with five distinct critical values. (The sign of a critical value is reversed
in the $\tau$ coordinate, since the equation is $\hat p(w)=-\tau$.)

\begin{lemma}[A connected cover generated by simple branches \src{05}]
\label{f6:kf:lem:transpositiongraph}
A complex polynomial of degree $n$ with $n-1$ simple critical points and
pairwise distinct critical values has monodromy group $S_n$.
\end{lemma}
\begin{proof}
This is the argument in the proof of Theorem~\ref{f6:thm:geometricS6}, which
uses nothing about the degree six. Remove the critical values from the target
and their preimages from the source; the result is an unramified
degree-$n$ cover whose source, the complex line with finitely many points
removed, is connected, so the monodromy is transitive. A small loop around
each critical value exchanges exactly two sheets, because the local map at
its unique critical point is analytically equivalent to $z\mapsto z^2$, and
such loops generate the fundamental group of the punctured target line. The
orbits of the group generated by these transpositions are the connected
components of their edge graph on the $n$ sheets; transitivity makes the
graph connected, and transpositions along the edges of a connected graph
generate all of $S_n$.
\end{proof}

\begin{proof}[Second proof of Theorem~\ref{f6:thm:geometricS6} \src{05}]
Away from the five values in \eqref{f6:kf:eq:criticalvalues}, the inverse
cover of the line \eqref{f6:kf:eq:line} is exactly $\hat p(w)+\tau=0$.
Lemma~\ref{f6:kf:lem:transpositiongraph} gives monodromy $S_6$ on that line.
Its image is a subgroup of the monodromy group over
$U=\{c\ne0,\ H\ne0\}$, which is in turn a subgroup of the permutations of six
sheets; it must therefore be $S_6$. Analytic continuation of algebraic root
germs preserves every polynomial relation over the base function field, so
it acts by automorphisms of the splitting field. Since the resulting
permutations already comprise $S_6$, the generic splitting-field group,
which embeds in $S_6$, is exactly $S_6$ as well.
\end{proof}

\begin{corollary}[No nontrivial intermediate inverse field \src{05}]
\label{f6:kf:cor:indecomposable}
The generic degree-six function-field extension of $\F$ has no proper
intermediate field. Consequently $\F$ cannot be written as a composition of
two dominant generically finite polynomial self-maps of $\A^5$, each of
geometric degree greater than one.
\end{corollary}
\begin{proof}
By Theorem~\ref{f6:kf:thm:reconstruction}, the source function field is
generated over the target function field by $w$. In the $S_6$ normal closure,
the field generated by one root is the fixed field of its point stabilizer
$S_5$. This is a maximal subgroup of $S_6$: any larger subgroup moves the
distinguished letter and, using its $S_5$ subgroup, is transitive on all six
letters; its point stabilizer is already $S_5$, so its order is
$6\cdot120=720$. Galois correspondence gives the field assertion. A
composition with both degrees greater than one would supply a proper
intermediate function field, contradicting it.
\end{proof}

This conclusion also says that a generic inverse cannot be described by
radicals over the rational function field of the target. The obstruction is
not the mere appearance of a sextic: some sextics are solvable by radicals.
It is the nonsolvable Galois group just established. The usual solvability
criterion and the elementary finite-field cycle method used below are
treated in \cite[Chapters~3--5]{Milne}.

\Needspace{5\baselineskip}
\subsection{An explicit rational certificate}
Consider the rational target
\[
 b_*=(1,0,0,1,0),\qquad
 f(w)=P_{b_*}(w)=2w^6-2w^5-w^3+w^2+1.
\]
Its discriminant is $-993580$. The following factorizations are in the
indicated finite fields:
\begin{align}
 f&\equiv2(w^6-w^5+w^3-w^2-1)&&\pmod3,\label{f6:eq:mod3}\\
 f&\equiv2(w-4)(w^5+3w^4-w^3+2w^2+2w-5)&&\pmod{13},\label{f6:eq:mod13}\\
 f&\equiv2(w+14)(w-16)(w-7)(w-5)(w^2+13w+14)&&\pmod{37}.
 \label{f6:eq:mod37}
\end{align}
The sextic in \eqref{f6:eq:mod3}, the quintic in \eqref{f6:eq:mod13}, and the
quadratic in \eqref{f6:eq:mod37} are irreducible in their respective fields.
The supplied verifier proves this with modular exponentiation and gcds,
not by accepting a black-box Galois-group output. Specifically, for a
monic polynomial $h$ of degree $n$ over $\mathbb F_p$, it checks
\[
 X^{p^n}\equiv X\pmod h,
 \qquad
 \gcd(h,X^{p^{n/\ell}}-X)=1
 \quad(\ell\text{ prime},\ \ell\mid n).
\]
These conditions are equivalent to irreducibility, by the factorization
of $X^{p^r}-X$ into the monic irreducibles whose degrees divide $r$.

\begin{theorem}[Arithmetic $S_6$ and inverse obstruction]
\label{f6:thm:arithmeticS6}
The polynomial $f$ is irreducible over $\Q$ and has Galois group $S_6$.
No one of its roots is expressible by radicals over $\Q$. In particular,
no full source point of $\F^{-1}(b_*)$ has all five coordinates expressible
by radicals over $\Q$.
\end{theorem}
\begin{proof}
Irreducibility modulo $3$ implies irreducibility over $\Q$, and hence a
transitive Galois action. At the three displayed primes the reductions
are squarefree and the leading coefficient remains a unit. The
factorization-cycle theorem gives a $5$-cycle from $13$ and a
transposition from $37$ \cite[Theorem 4.28]{Milne}. To fit the monic
integer formulation exactly, apply it to $2^5f(W/2)$; at these odd
primes the change of variable preserves factor degrees.

A transitive subgroup of $S_6$ containing a $5$-cycle has a point
stabilizer transitive on the other five labels, and is therefore
$2$-transitive. Conjugating the transposition gives every transposition.
Hence the group is $S_6$, in agreement with the standard group criterion
\cite[Lemma 4.31]{Milne}.

An irreducible polynomial with nonsolvable splitting group cannot have
a root expressible by radicals: all conjugates of such a root would
belong to a common solvable normal radical extension after adjoining
suitable roots of unity, contradicting nonsolvability of $S_6$
\cite[\S5, solvability theorem]{Milne}.
Finally $w=\gamma u_1$ is a polynomial in the source coordinates with
rational coefficients. If an entire inverse vector had radical
coordinates, its corresponding root $w$ would be radical as well.
\end{proof}

This result does not forbid an exact inverse expressed by algebraic
numbers or by a root-isolation representation. It forbids a radical-only
inverse vector in this explicit fiber. It is therefore directly relevant
to the repository's distinction between exact algebraic solving and
radical-solvability certificates.

\Needspace{5\baselineskip}
\subsection{A second rational target with six nonradical inverse points
\src{05}}
\label{f6:kf:sec:arithmetic}
Source~05 certifies a different target, by different primes and a different
group-theoretic deduction. Consider
\begin{equation}\label{f6:kf:eq:qstar}
b_{**}=(1,-1,0,1,1),\qquad
R_{**}(w)=P_{b_{**}}(w)=2w^6-2w^5-w^3+w+1,
\end{equation}
with $\Disc R_{**}=-1513772=-4\cdot13\cdot43\cdot677$. It is squarefree, so
the fiber contains six complex points; its effective parameter is
$(a,b,d)=(0,1,1)$, and it lies on the Morse line \eqref{f6:kf:eq:line} at
$\tau=1$.

\begin{theorem}[Arithmetic $S_6$ certificate \src{05}]
\label{f6:kf:prop:arithmetic}
The polynomial $R_{**}$ is irreducible over $\Q$ and has Galois group $S_6$.
In particular, $b_{**}$ has no rational preimage, and none of its six source
points has all coordinates expressible by radicals over $\Q$.
\end{theorem}
\begin{proof}
The following factorizations hold in the indicated finite fields:
\begin{align}
R_{**}&\equiv2(w^6-w^5+3w^3-3w-3)&&\pmod7,\label{f6:kf:eq:mod7}\\
R_{**}&\equiv2(w+3)(w^5-4w^4+w^3+2w^2+5w+2)&&\pmod{11},\label{f6:kf:eq:mod11}\\
R_{**}&\equiv2(w+76)(w+98)(w-132)(w-114)(w^2+71w+134)&&\pmod{269}.
\label{f6:kf:eq:mod269}
\end{align}
The degree-six, degree-five and degree-two factors, respectively, are
irreducible; the Frobenius tests proving this are listed below, and the
supplied script of source~05 executes them with modular polynomial
arithmetic. The three primes divide neither the leading coefficient nor the
discriminant.

Reduction modulo $7$ proves irreducibility over $\Q$, so the Galois group is
transitive. For a squarefree reduction at a good prime, Frobenius permutes
the roots in cycles whose lengths are the degrees of the irreducible factors;
lifting Frobenius to a decomposition group gives an element of the
characteristic-zero Galois group with the same cycle type. The reductions
therefore supply a $6$-cycle, a $5$-cycle with one fixed point, and a
transposition with four fixed points.

A transitive group of degree six containing a $5$-cycle is primitive. Indeed,
a nontrivial block system would have blocks of size two or three. An element
of order five must fix each block, since there are at most three blocks, and
must act trivially within each block, since each has size at most three;
this contradicts the presence of the $5$-cycle. Finally, a primitive
permutation group containing a transposition is the full symmetric group. To
see this directly, take all conjugates of that transposition and form their
edge graph. The connected components are blocks for the original group, so
primitivity forces the graph to be connected, and the edge transpositions
then generate $S_6$.

A rational source point would yield a rational root $w=\gamma u_1$, contrary
to irreducibility. If all coordinates of a source point were expressible by
radicals, so would $w$ be, because $\gamma$ and $u_1$ are polynomials in the
source coordinates with rational coefficients. An irreducible polynomial with
Galois group $S_6$ has no root expressible by radicals: the normal closure of
a radical expression has solvable Galois group after adjoining roots of unity,
whereas the normal closure of any one root of $R_{**}$ is its $S_6$ splitting
field.
\end{proof}

The Frobenius tests are those displayed before
Theorem~\ref{f6:thm:arithmeticS6}. For a monic $f\in\mathbb F_p[w]$ of degree
$n$, $w^{p^n}\equiv w\pmod f$ together with $\gcd(f,w^{p^{n/\ell}}-w)=1$ for
every prime $\ell\mid n$ is sufficient and necessary for irreducibility:
$w^{p^n}-w$ is the product of the distinct monic irreducibles whose degrees
divide $n$, the gcd tests remove factors of every proper divisor degree
(every proper divisor of $n$ divides some $n/\ell$), and the congruence also
forces squarefreeness. For the three nonlinear factors above the complete
tests are:
\begin{center}
\renewcommand{\arraystretch}{1.25}
\begin{tabular}{@{}ccl@{}}
\toprule
Field & Degree & Tests\\
\midrule
$\mathbb F_7$ & $6$ & $w^{7^6}=w\pmod f$; gcds for $7^3$ and $7^2$ equal $1$\\
$\mathbb F_{11}$ & $5$ & $w^{11^5}=w\pmod f$; gcd for $11$ equals $1$\\
$\mathbb F_{269}$ & $2$ & $w^{269^2}=w\pmod f$; gcd for $269$ equals $1$\\
\bottomrule
\end{tabular}
\end{center}
All congruences are evaluated by repeated squaring modulo a polynomial of
degree at most six; source~05's program does not expand polynomials of degree
$7^6$ or $11^5$ in the ordinary polynomial ring. It also checks each displayed
factorization and that the primes are good for the discriminant. The
Galois-group certificate is thus small, reproducible, and independent of a
black-box command returning a group name.

The two sources' certificates are independent: source~02 uses
$b_*=(1,0,0,1,0)$ with primes $3,13,37$ and a $2$-transitivity argument;
source~05 uses $b_{**}$ with primes $7,11,269$ and a primitivity argument.
The conclusions for the omitted point \eqref{f6:kf:eq:rationalmissing} and
for $b_*$, $b_{**}$ are deliberately different. The former target has no
geometric preimage. The latter have six geometric preimages, none rational,
with a proven radical obstruction. This distinction prevents an arithmetic
non-surjectivity statement from being mistaken for a geometric one.

\Needspace{6\baselineskip}
\section{Real fibers and arithmetic inversion}
\label{f6:sec:real}

\begin{theorem}[Real cardinalities]\label{f6:thm:real}
For the real polynomial map $\F:\R^5\to\R^5$, the possible fiber
cardinalities are exactly $\{0,1,2,3,4\}$. If $c\ne0$, the possible
values are contained in $\{0,1,2,4\}$. If $c=0$, the count is $3$ for
$A^2+4B>0$ and $1$ for $A^2+4B\le0$.
\end{theorem}
\begin{proof}
For $c\ne0$ the real source points correspond to the real simple roots
of $P_b$. Suppose all six roots $r_i$ were real, counted with
multiplicity. Newton's identities applied to \eqref{f6:eq:fixedcoeff} give
\[
 \sum_i r_i=1,\qquad \sum_i r_i^2=1,\qquad
 \sum_i r_i^3=\frac52.
\]
But
\[
 \left|\sum_i r_i^3\right|\le\sum_i|r_i|^3
 \le\left(\sum_i r_i^2\right)^{3/2}=1,
\]
a contradiction. Five simple real roots would force the remaining root
to be real as well. Exactly three simple real roots would leave a
real cubic factor with no additional simple real root. Its odd degree
forces a real root, and that root must be triple; otherwise a further
simple real root remains. This would again make all six roots real.
Thus only $0,1,2,4$ simple real roots can occur.

For $c=0$, use $q_b=v^2+Av-B$ and the one-point section. Two distinct
real roots give three points; a double root is deleted, and nonreal
roots give no real points. This proves the stated criterion.

All five global values occur. Zero occurs at $b_*$: write
$p(w)=w^2(w-1)(2w^3-1)$. It is nonnegative outside
$(2^{-1/3},1)$, while on that interval $-1<p(w)<0$. Thus $p(w)+1>0$
on all of $\R$. One occurs at the origin, two at $(1,0,0,0,0)$, and
three at $(0,0,1,0,0)$. Four occurs at $(1,-8,0,-6,70)$, whose sextic is
\[
 2w^6-2w^5-w^3-7w^2+14w-6.
\]
It has the simple root $1$, and sign changes on
$(-2,-1)$, $(0,9/10)$, and $(11/10,2)$. These give four distinct real
roots; the preceding bound proves there are no others. Exact Sturm
counts independently verify all five examples.
\end{proof}

A particularly instructive target is therefore $b_*$. Its geometric
fiber has six points, its real fiber is empty, and its rational fiber is
empty. These statements coexist without contradiction: geometric
surjectivity onto $\A^5\setminus M$ does not assert real or rational
surjectivity.

\Needspace{5\baselineskip}
\subsection{An exact inversion procedure over a computable field}
The results give a compact algorithm which never expands or eliminates
all $2180$ monomials of the original map.
\begin{enumerate}[label=\textbf{Step \arabic*.}]
\item Form $\kk,T,L$ and $q_b$ by \eqref{f6:eq:TL}--\eqref{f6:eq:q}. For geometric
cardinality, compute the gcds in \eqref{f6:eq:simplepart} and take the degree,
adding one when $c=0$.
\item For inverse points, factor or isolate the roots of $\simp(q_b)$ in
the desired field. Over $\Q$, degree-one factors are exactly the rational
roots; higher-degree factors represent algebraic inverse points.
\item At each simple root evaluate \eqref{f6:eq:inversechart},
\eqref{f6:eq:stageinverse}, and \eqref{f6:eq:sourceinverse}. The only variable
denominator is a power of $q_b'(v)$, whose nonvanishing has already been
certified.
\item If $c=0$, also solve the triangular equations \eqref{f6:eq:section}.
This supplies the unique point not represented by the $x\ne0$ chart.
\end{enumerate}
Polynomial degree is at most six throughout the root computation, and
drops to two on $c=0$. The formula handles multiplicities exactly; a
near-zero numerical derivative must never be substituted for a gcd test.
The accompanying Python module exposes the root polynomial, simple-part
calculation, rational reconstruction, and section inverse as reusable
functions.

\Needspace{6\baselineskip}
\section{Verification, proof boundaries, and formalization}
\label{f6:sec:verification}

\Needspace{5\baselineskip}
\subsection{What the reproducibility package checks}
Source~02's package contains a sparse rational construction of $\F$, a separate
verification driver, two coefficient files, and a saved verification log.
The driver completed \textbf{63 exact checks} under Python 3.13.5 and
SymPy 1.14.0. No floating-point arithmetic or network access is used.

The checks include cancellation of the apparent denominators, degree and
term counts, the entire $x=0$ section, the stage inverse, the stage and
sweep determinants, the derivative identity, the rescaled root equation,
the three normalization relations, the exceptional quadratic, the exact
$c^2$ discriminant factor and its leading coefficient, exclusion of the
other all-multiple partitions, every cardinality witness, exact inverse
round trips including $\gamma=0$, the Morse discriminant, finite-field
irreducibility certificates, and real Sturm counts.

The saved log is a reproducibility record, not an independent proof
assistant certificate. In particular, a finite selection of inverse
round trips cannot establish a chart isomorphism. That is why
Theorem~\ref{f6:thm:chart} proves the identities on a dense open set and
then explicitly proves that they extend over the exceptional locus.
Likewise, normality, openness of the embedding, and the monodromy
argument are proved mathematically; they are not inferred from a
computer-algebra ``PASS'' message.

\Needspace{5\baselineskip}
\subsection{The exact checks of source~05 \src{05}}
Source~05's companion program is independent of source~02's. It constructs
the sweep, rather than entering a precomputed large map. It checks the sweep
Jacobian, both compositions of the stage with its inverse, and the stage
determinant. It then forms the four quotient polynomials $\widehat S_i$ and
proves the required divisibilities by collecting all coefficients that could
obstruct them; the boundary jets and their triangular Jacobian are checked
exactly. For the fiber analysis it checks the sextic elimination, identity
\eqref{f6:eq:derivative}, and all four sweep equations modulo $P$ after
reconstruction. It verifies the square polynomial and its discriminant, the
sample factorization table, the triple-root cofactor and its discriminant,
the omitted-surface parametrization, and the $c$-adic discriminant
coefficient. Finally it checks the critical-value polynomial of the second
Morse slice and the three modular factorizations, with irreducibility tests
performed by repeated squaring and gcd computations. The recorded run
(Python 3.13.5, SymPy 1.14.0) reports twelve PASS groups, then
\texttt{ALL CHECKS PASSED}, followed by two lines restating the proof
boundary.

All these are rational or finite-field calculations. No floating-point
roots, random samples, heuristic factor selections or remote computations
enter the certificate. The full coefficient expansion of $\Delta$ is
regenerated by the program. The script is an independent computational audit
of explicit identities, not an alternate proof of every theorem: it cannot
replace the proof that the seven multiplicity partitions are exhaustive, the
local analytic coordinate argument, the properness argument, or the
group-theoretic deductions. These are given in the text.

Neither suite checks the trace-form statement of Theorem~\ref{f6:thm:H}
($\Disc(\mathcal B)=H/256$), nor source~05's local normal forms beyond the
quartic discriminant \eqref{f6:kf:eq:quarticdisc}; the intake's independent
spot check (not shipped) confirmed the trace-form identity at three rational
targets. In preparing this merge,
the renamed formulas of source~05 printed here (the quotient polynomials of
Appendix~\ref{f6:kf:app:polynomials}, identity \eqref{f6:kf:eq:sweepdet},
the restriction \eqref{f6:kf:eq:Egamma0}, the triple-root family, the second
Morse slice, and the factorizations modulo $7,11,269$) were re-checked with
SymPy~1.14.0 against the construction \eqref{f6:eq:stage}--\eqref{f6:eq:map}.
That check is not shipped.

\Needspace{5\baselineskip}
\subsection{A suitable Lean development boundary}
A formalization should begin with the small algebra, not with a giant
five-variable graph ideal. The natural dependency order is:
\begin{enumerate}[label=\textup{(\arabic*)}]
\item Prove the elimination and derivative identities in a polynomial
ring. These are finite ring equalities with rational coefficients.
\item Formalize the multiplicity-one factor construction
\eqref{f6:eq:simplepart}, including the zero-characteristic hypotheses and
constant-polynomial edge case. This yields a reusable exact fiber-count
checker.
\item Establish the six-element algebra and its multiplication rules.
Prove spanning by closure under the two generators and independence by
the degree-six field relation. Then derive the discriminant basis change.
\item Formalize the source chart and its exceptional section, either as
localized coordinate-ring isomorphisms or first as field-valued point
bijections with explicit nonvanishing hypotheses.
\item Add the geometric layer: normality of the three charts, the open
immersion criterion, and the branch/nonproperness interpretation. This
stage needs substantially more algebraic-geometric infrastructure than
the finite polynomial certificates.
\end{enumerate}
No Lean source is supplied as if it were checked. A future kernel audit
should report separately which finite identities, pointwise statements,
and scheme-theoretic conclusions have actually been formalized.

\Needspace{5\baselineskip}
\subsection{A narrow, layered Lean route \src{05}}
Source~05 proposes a parallel plan, also explicitly a proposed implementation
plan and not a report of completed Lean work. Its second layer corresponds to
steps~(1) and~(4) above, its third to the exceptional section of step~(4),
its fourth to step~(2) together with the Galois certificates, and its fifth
to step~(5). Its first layer (the multiplicity patterns) has no counterpart
above, and it has no counterpart of step~(3). The repository's style of connecting exact algebraic
certificates to mathematical semantics suggests the following order.

\paragraph{First layer: the sextic family.} Define
$P=2w^6-2w^5-w^3+aw^2+bw+d$ as a univariate polynomial over a
characteristic-zero field. Prove the coefficient identities and the
power-sum constraints. Formalize the two-root contradiction
(Lemma~\ref{f6:kf:lem:threeroots}) separately from any multivariate
geometry. The all-multiple-root theorem can then be proved by multiplicity
partitions and coefficient comparison. This yields the essential
omitted-parameter result in a small algebraic development.

\paragraph{Second layer: reconstruction.} Represent the stage and the sweep as
multivariate polynomials, and use normalized polynomial certificates for
\eqref{f6:eq:derivative}. State the reconstruction theorem first for points
with explicit nonzero hypotheses $c\ne0$ and $P_b'(w)\ne0$. Avoid a
definition that treats field division at zero as if it were a removable
quotient. The polynomial extension (Proposition~\ref{f6:prop:keller}) must be
built from divisibility, with the quotient identities proved before
evaluation.

\paragraph{Third layer: the hyperplane.} The triangular map
\eqref{f6:eq:section} and the equation $x^2(A^2+4B)=1$ require no scheme
theory. They can provide a complete set-theoretic fiber-count theorem,
including the exceptional locus, before any formalization of properness or
analytic monodromy is attempted.

\paragraph{Fourth layer: finite algebra and certificates.} Formalize the
simple-root-factor formula and its interpretation. The finite-field
irreducibility checks can be encoded as polynomial remainder and gcd
certificates. Connecting the cycle types to $S_6$ requires an explicit
Frobenius bridge and elementary permutation-group lemmas; it should not be
hidden behind an unchecked external Galois-group command.

\paragraph{Fifth layer: geometry and topology.} Scheme localization, the
normalization of $\dd$, analytic preparation, and the covering-space argument
are substantial separate interfaces. A formalized algebraic fiber theorem
should not be advertised as a formalized nonproperness or monodromy theorem
until those interfaces are completed.

This layered plan gives useful, independently reviewable endpoints. It also
keeps the small exact certificates outside the trusted boundary until a
proved checker or a direct kernel proof verifies their semantic meaning.
None of the layers exists in ProveIt at the time of this report
(Section~\ref{f6:sec:merge}).

\Needspace{5\baselineskip}
\subsection{Why the argument is not a numerical discovery claim}
The important transitions are algebraic. A root is deleted because an
exact derivative is zero. An entire fiber is empty because an exhaustive
list of multiplicity partitions reduces to one square. The finite model
is normal because its charts are smooth, not because sample fibers look
correct. The Galois group is $S_6$ because of exact permutations forced
by finite-field factorizations and a connected Morse cover. These
arguments continue to hold at targets which numerical root solvers
would find highly ill-conditioned.

\Needspace{6\baselineskip}
\section{Further research questions and proposed projects}
\label{f6:sec:future}
The following questions are specific continuations of the results above.
They are proposed directions, not additional theorems established here.

\Needspace{5\baselineskip}
\subsection{Complete the adjacent higher-dimensional fiber problems}
\textbf{1. The degree-twelve map $F_7$.} Gao supplies a univariate
fiber polynomial of degree twelve \cite[\S4.6]{Gao}. Determine whether
its exceptional degeneration is resolved by successive versions of
Theorem~\ref{f6:thm:general}. A concrete first deliverable is an integral
basis and a factorization of its normalized discriminant, followed by
an exhaustive empty-fiber classification. The $F_6$ analysis suggests
that the exceptional hyperplane must be treated independently of the
generic elimination equation.

\textbf{2. The maps $F_4$ and $F_5$.} Their published quintic tangency
and degree-ten resultant descriptions leave related fiber problems
\cite[Appendix A.3]{Gao}. Derive a derivative or subresultant criterion
that distinguishes an actual source point from a deleted tangency
parameter. The main risk is that a resultant introduces extraneous
solutions; each proposed root representation should be accompanied by
an explicit inverse chart.

\Needspace{5\baselineskip}
\subsection{Refine the geometry of the sextic example}
\textbf{3. Singular strata of $H$.} The simple-root gcd criterion gives
the correct count at every point, but does not constitute a Whitney
stratification. Determine the singular locus of the discriminant via
\eqref{f6:eq:incidence}, separating triple roots, two double roots, and
higher coincidences. Describe their intersections with $c=0$ and the
omitted surface $M$. An attractive deliverable is a small set of
parametrizations or subresultant ideals rather than a large expanded
elimination basis.

\emph{Status after the merge.} Partly answered by source~05, and re-scoped.
Theorem~\ref{f6:kf:thm:localtypes} gives the local analytic type of $\dd$ at
every point of coefficient space and the singular-locus criterion (a root of
multiplicity at least three, or at least two distinct repeated roots);
Proposition~\ref{f6:kf:prop:Dnorm} gives the normalization of $\dd$;
Corollary~\ref{f6:kf:cor:crossings} describes $V(H)$ along $M$; and
Theorem~\ref{f6:thm:nonproper} gives its behaviour along $c=0$. What remains
open is a Whitney stratification, explicit parametrizations or ideals of the
singular strata of $V(H)$ in the five-dimensional target, and their full
intersection pattern with $c=0$. Source~05 likewise claims no Whitney
stratification.

\textbf{4. Quantitative escape rates.} For each deleted component $E_i$,
construct explicit Laurent-series source arcs whose images tend to a
prescribed point of $c=0$. Compute the leading pole orders of all five
source coordinates. Repeat near a generic ramification point and near
the triple-double locus $M$. This would distinguish the unramified and
ramified mechanisms dynamically, not just set-theoretically.

\textbf{5. The smooth affine normalization $Z$.} Determine its divisor
class group, Picard group, and algebraic automorphisms. Is the explicit
six-generator order monogenic after a target automorphism or after a
smaller localization? Any answer must respect the quadratic-plus-quartic
special fiber \eqref{f6:eq:CRT}; the naive power basis cannot do so.

\Needspace{5\baselineskip}
\subsection{Extend the algebraic and arithmetic methods}
\textbf{6. Higher contractions.} Replace $w^2J-cTw-c^2L$ by a family
whose leading degeneration contains $w^r$, with $r\ge3$. Find hypotheses
under which $r-1$ explicit divided basis elements suffice to normalize
it. The desired theorem should state a finite-free basis, multiplication
rules, discriminant correction, and checkable normality conditions.

\textbf{7. Arithmetic image and solvable fibers.} Classify rational
parameter subvarieties on which the sextic Galois group drops to a
solvable subgroup. Distinguish those with a rational source point from
those with only nonrational radical inverse points. The explicit $S_6$
certificate is a starting obstruction, not a classification of all
specializations. This problem connects directly to ProveIt's
radical-solvability and exact polynomial-solving projects.

\textbf{8. Good reduction.} Analyze the root chart and normalization in
positive characteristic, separating primes excluded by the map's
coefficients or determinant from primes affecting discriminants such as
$401$. The characteristic-zero proof must not simply be reduced modulo
an arbitrary prime: derivative multiplicities, separability, and the
stage determinant require independent checks.

\Needspace{5\baselineskip}
\subsection{Formal and structural questions}
\textbf{9. A certified root-chart toolkit.} Build a reusable library that
takes a polynomial identity for a graph, a reconstruction tuple, and
nonvanishing certificates, and returns a checked fiber bijection. Add
support for exceptional source components instead of silently localizing
them away. This would serve the Jacobian, radical-expression, and
algebraic-number parts of ProveIt simultaneously.

\textbf{10. Rigidity versus finite-cover complexity.} Gao's middle-branch
rigidity questions concern when a higher-dimensional construction is
only a suspension of a lower-dimensional one \cite[\S4.5.2]{Gao}.
Investigate whether normalized-cover invariants, especially monodromy
and the number of unramified boundary components, obstruct such
suspensions. The present $S_6$ computation alone is not a proof of a
new rigidity theorem; a comparison theorem for finite normalizations
would be needed. Source~05's Corollary~\ref{f6:kf:cor:nostabilization}
excludes a suspension of $F_6$ from dimension at most three, by the
geometry of the omitted surface; direction~13 below asks about dimension
four.

\Needspace{5\baselineskip}
\subsection{Directions from source~05 \src{05}}
\label{f6:kf:sec:future}
Source~05 lists eight further questions, numbered 11--18 here. The completed
calculations make them accessible; they are not assertions that the
corresponding results already hold. Where source~02's results answer or
overlap one, this is stated.

\textbf{11. A global finite completion across $c=0$ (answered by
source~02).} Over $c\ne0$, \eqref{f6:kf:eq:completion} gives a smooth finite
completion with a transparent deleted ramification divisor. Source~05 asked
for an explicit finite normalization over the whole target $\A^5$, with the
original affine source located inside it; in particular, for the boundary
components above $c=0$, their degrees, and how the three lost sheets interact
with the two-sheet discriminant loss at \eqref{f6:kf:eq:boundaryintersection}.
Theorems~\ref{f6:thm:normalization} and~\ref{f6:thm:boundary} answer this:
$Z=\Spec\normal$ is finite free of rank six and smooth; the source is
$Z\setminus(\Ram(\pi)\cup E_1\cup E_2\cup E_3)$; each $E_i\simeq\A^4$ maps
isomorphically (degree one) onto $c=0$, is unramified and is disjoint from
$\Ram(\pi)$; and by \eqref{f6:eq:CRT} the two remaining sheets over $c=0$
form the quadratic factor $\zeta^2+A\zeta-B$, which ramifies exactly along
$A^2+4B=0$. Refinements, such as escape rates along each component, are
source~02's question~4.

\textbf{12. Higher-codimension omitted sets.} The omitted surface arises from
three independent double roots of a degree-six fiber polynomial. Can a
tangent-sweep construction realize $r$ independent double roots, with an
omitted locus of codimension $r$, for arbitrary $r$? Within a local model
consisting of disjoint simple folds and no surviving sheets, degree at least
$2r$ is necessary simply by counting the two sheets of each fold; the present
example realizes equality for $r=3$. A theorem for general Keller maps would
require much more than this local counting observation. A practical route is
to start with a family of degree-$2r$ polynomials having a unique square
member and an $r$-dimensional transversal coefficient slice, then solve the
stage divisibilities that turn its ramification into escape. The difficult
part is not the square polynomial, but constructing a global polynomial
source with constant Jacobian.

\textbf{13. The remaining four-dimensional stabilization question.}
Corollary~\ref{f6:kf:cor:nostabilization} excludes every stabilization from
dimension at most three. Is $F_6$ nevertheless equivalent to
$\Phi\times\id_{\A^1}$ for a four-dimensional Keller map $\Phi$ of geometric
degree six? Its omitted curve would have to have a cylinder isomorphic to
$\Gm\times\A^1$, so the omitted surface is compatible with such a possibility
rather than excluding it. The embedding of that surface, the normalization of
the escape divisor, and the distribution of boundary sheets in a global
finite completion may decide this remaining case. (That completion is now
$Z$; see direction~11.) The composition obstruction of
Corollary~\ref{f6:kf:cor:indecomposable} does not settle it: adding an
identity coordinate does not factor the generic degree into smaller degrees.

\textbf{14. The degree-twelve fiber polynomial of $F_7$ (overlaps
question~1).} Apply the derivative-localization method to the other
five-dimensional construction in \cite{Gao}. A useful first target is not the
entire discriminant expansion, but a characterization of all fiber
polynomials without simple roots. Degree twelve admits many more multiplicity
partitions, so the first questions are which Newton-sum constraints survive
and whether a small effective parameter space remains transverse to every
allowed contact pattern.

\textbf{15. Real image and real fiber chambers (re-scoped).} The
reconstruction formulas are defined over $\Q$, so for real targets with
$c\ne0$ the real fiber is exactly the set of real simple roots of $P_b$; on
$c=0$ the real count is three when $A^2+4B>0$ and one otherwise.
Theorem~\ref{f6:thm:real} of source~02 has since determined the possible real
fiber cardinalities, $\{0,1,2,3,4\}$, with $\{0,1,2,4\}$ containing the values
off $c=0$. What remains open is source~05's actual request: the connected
real chambers, the real fiber count on each, and the semialgebraic boundary
of the real image, and whether each value in $\{0,1,2,4\}$ is attained off
$c=0$. Sturm sequences for the exact simple-root factor provide an immediate
computational route, but a useful result should describe the chambers
geometrically rather than merely output a quantifier-elimination formula.

\textbf{16. Arithmetic specialization and positive characteristic (overlaps
questions~7 and~8).} For rational targets, determine the distribution of
specialized Galois groups, the arithmetic of one-point and two-point fibers,
and the loci admitting rational source points. The explicit $S_6$ targets
give starting certificates, not such a distribution theorem. The
construction also invites reduction modulo primes, but characteristic-zero
conclusions must not be transferred indiscriminately. The denominator $5$,
the determinant $-290$, the cubic discriminant $-401/128$, and the integer
moment obstructions already reveal several exceptional primes. Determine
exactly which fiber patterns and omitted-locus formulas survive at each
prime, and separate inseparability from genuine changes in the coefficient
constraints.

\textbf{17. The middle-branch rigidity problem remains separate (overlaps
question~10).} Gao's Problem~4.8 concerns whether the entire middle branch of
a mixed direction-field construction is a suspension. The present results
analyze a bottom-branch example and do not answer that question. A
productive next step is to derive a fiber-polynomial or incidence invariant
for the middle branch and test whether the suspension restriction imposes a
detectable factorization on its monodromy or ramification geometry. Such a
test could produce either a rigidity proof or a systematically searchable
counterexample class.

\textbf{18. Formal, reusable fiber-classification infrastructure (overlaps
question~9).} A reusable library should take a monic polynomial family and
an explicit localization condition, compute the simple-root factor, and
connect it to the fibers of a polynomial map via a proved reconstruction
equivalence. The present case is small enough to serve as a complete
benchmark: it includes generic fibers, repeated roots, an exceptional
hyperplane, an omitted locus, and a finite-field Galois certificate.
Formalizing this example would be more valuable than checking isolated
collisions, because it would verify the transition from elimination algebra
to a global classification.

\Needspace{6\baselineskip}
\section{Conclusion}
\label{f6:sec:conclusion}
The central open task selected from the repository's surrounding research
program was not the existence of another collision, but the exact
behavior of every fiber of an already explicit high-degree map.
The derivative identity makes that task accessible. The uniform chart
then exposes the exceptional quadratic, and one divided basis element
repairs the finite model. These steps determine the image, all geometric
and real cardinalities, the branch and nonproperness loci, and a concrete
radical obstruction.

The resulting model is small: a sextic, a quartic, a quadratic exceptional
fiber, and three multiplication rules. It retains the geometry hidden by
the large expanded map and provides a practical target for independent
review, further generalization, and formal verification.

Source~05 draws the same conclusion from the other side \src{05}. The central
fact is not simply that a sextic appears in an inverse problem: its
derivative is exactly the multiplier that decides which inverse points remain
in affine space. Once this identity is proved, the question of omitted
targets becomes a question about a rigid coefficient slice with no simple
roots, and Newton sums reduce it to a single square of a cubic. The
exceptional hyperplane contributes no omitted points and is controlled by one
quadratic equation. The resulting picture is exact: a smooth
codimension-three omitted surface, all six possible fiber sizes and all seven
multiplicity patterns, a two-component nonproperness divisor, the
normalization and all local analytic types of the discriminant, a
dimension-three stabilization obstruction, and full symmetric monodromy.
These are concrete extensions of the polynomial-map research direction
represented in ProveIt. Their algebraic core is accompanied by exact
certificates; the remaining formalization and research problems are
separated from the proved statements.

\appendix
\Needspace{6\baselineskip}
\section{Reproducing the calculations}
\label{f6:app:reproduce}
Both sources' verification suites are shipped, byte-identical to the
delivery but renamed with the prefixes \code{02-six-sheets-} and
\code{05-keller-fibers-}. Each script writes its outputs \emph{beside itself
under its delivered, unprefixed names}, and source~02's verifier imports
\code{build_map} by its delivered module name. Rerun them therefore only on a
copy laid out as delivered, never in the report directory; this keeps the
shipped records intact.

\emph{Source~02} (SymPy~1.14.0, pinned in
\code{data/02-six-sheets-requirements.txt}). In a scratch directory with
subdirectories \code{code} and \code{data}, copy
\code{code/02-six-sheets-build_map.py} to \code{code/build_map.py} and
\code{code/02-six-sheets-verify.py} to \code{code/verify.py}, then run
\begin{lstlisting}
python -m pip install sympy==1.14.0
python code/verify.py
\end{lstlisting}
from the scratch directory. It rewrites \code{data/map_coefficients.json},
\code{data/normalized_discriminant.json} and
\code{data/verification_report.txt} there; compare them with the shipped
\code{data/02-six-sheets-*} files. Run standalone, \code{build_map.py} also
writes \code{data/zero_section.txt}, which is not shipped.

\emph{Source~05} (same pin, \code{data/05-keller-fibers-requirements.txt}).
Copy \code{code/05-keller-fibers-verify.py} to \code{verify.py} in a scratch
directory and run
\begin{lstlisting}
python verify.py > verification.txt
\end{lstlisting}
It writes \code{sextic_discriminant.txt} beside itself even without
redirection; compare the two outputs with
\code{data/05-keller-fibers-verification.txt} and
\code{data/05-keller-fibers-sextic_discriminant.txt}. The function
\code{complex_fiber_size}, imported as \code{from verify import
complex_fiber_size}, accepts an exact rational target $(c,A,B,D,E)$ and
counts geometric points by the simple-root factor and
\eqref{f6:eq:c0count}; it is not a count of rational or real preimages.

On Windows, text-mode writes add carriage returns; the shipped files have
none, and the intake's reruns matched them byte for byte after stripping
carriage returns. Neither suite needs network access after installation, and
neither needs a \LaTeX\ installation. This report compiles independently of
Python with
\begin{lstlisting}
latexmk -pdf -interaction=nonstopmode -halt-on-error article.tex
\end{lstlisting}

\begin{center}\small
\renewcommand{\arraystretch}{1.25}
\begin{tabularx}{\textwidth}{@{}>{\raggedright\arraybackslash}p{0.47\textwidth}X@{}}
\toprule
Shipped file (delivered name) & Purpose\\
\midrule
\code{code/02-six-sheets-build_map.py} (\code{code/build_map.py}) & Source~02: sparse rational construction and exact denominator cancellation.\\
\code{code/02-six-sheets-verify.py} (\code{code/verify.py}) & Source~02: the 63-check certificate and reusable rational inverse functions.\\
\code{data/02-six-sheets-map_coefficients.json} & All five polynomial coordinates; monomial exponent order is $(x,y,z_1,z_2,z_3)$.\\
\code{data/02-six-sheets-normalized_discriminant.json} & Sparse $H$ in coordinate order $(c,A,B,D,\kk)$; $\kk=E+B^2-AD$ (written \code{K} in the file).\\
\code{data/02-six-sheets-verification_report.txt} & Saved output of source~02's completed exact run.\\
\code{data/02-six-sheets-requirements.txt} & \code{sympy==1.14.0}.\\
\code{02-six-sheets-SOURCE_AUDIT.md} & Source~02: repository revision, source scope, and novelty boundary.\\
\code{code/05-keller-fibers-verify.py} (\code{verify.py}) & Source~05: exact rational and finite-field checks, and \code{complex_fiber_size}.\\
\code{data/05-keller-fibers-verification.txt} & Saved output of source~05's run (twelve PASS groups).\\
\code{data/05-keller-fibers-sextic_discriminant.txt} & The expanded $54$-term $\Delta(a,b,d)$, written in source~05's coefficient names $(a,b,c)$, so its $c$ is $d$ here.\\
\code{data/05-keller-fibers-requirements.txt} & \code{sympy==1.14.0}.\\
\code{05-keller-fibers-SOURCE_AUDIT.md} & Source~05: pinned sources, comparison scope, and attribution.\\
\bottomrule
\end{tabularx}
\end{center}

Every coefficient in source~02's JSON files is a rational string; no
precision parameter is needed. Its verifier regenerates the coefficient files
from the formulas rather than accepting them as trusted input. A failure
raises an exception and prevents the final all-checks-passed message. There
are no random checks in either verifier.

\Needspace{6\baselineskip}
\section{Source audit and attribution}
\label{f6:app:sources}
The repository reference observed during the audit was
\begin{center}\small
\code{e21766d04c2b8a9b2cdba0cd43e563024b3bd1b9}.
\end{center}
The main directly inspected project files were
\begin{center}\small
\code{Algebra/JacobianConjecture/README.md}\\
\code{Algebra/JacobianConjecture/Research/README.md}.
\end{center}
Their returned blob hashes were respectively
\begin{center}\small
\code{67e25dcb94cb7a38fb508ea65e3a35455df859ac}\\
\code{f0605a788b98bed23df91318aa2069e6802fc66e}.
\end{center}
The root README and a walkthrough search result were used for orientation.
The root README's general formalization claims were not independently
rebuilt in this work, and none is used as a substitute for the specific
algebraic proofs above.

The main external mathematical source was Gao's version 1, dated
31 July 2026, particularly \S4.5.1, Theorem 4.5, the concluding remarks,
and Appendix A.3. Its definition of $F_6$, determinant, generic degree,
and elimination polynomial are explicitly credited. The derivative
identity, uniform chart, complete cardinality criterion, exact omitted
surface, six-element normalization, normalized discriminant and boundary,
and the geometric and arithmetic $S_6$ proofs are the extensions developed
in this article. This describes their relation to the inspected sources,
not a guarantee that no equivalent result exists elsewhere.

The primary literature search checked the cited paper's available version
and searches combining its identifier or author with $F_6$, fiber,
normalization, and Jacobian terminology. No directly matching extension
was identified in that bounded search. It was not an exhaustive review
of all unpublished work, repository branches, or later inaccessible
material. The Stacks Project and Milne are cited for the standard
geometric and Galois-theoretic facts used in the proofs.

\emph{Source~05's audit} \src{05}. Source~05 inspected the same pinned commit.
Its principal repository sources were
\code{Algebra/JacobianConjecture/Lean/JacobianConjecture/Counterexample.lean},
\code{Algebra/JacobianConjecture/Research/README.md}, and
\code{ProveIt_Walkthrough.tex}, specifically the walkthrough's reference to
Gao's account as external context for the repository's formal determinant
and collision identities \cite{ProveItWalkthrough}. It also inspected the
Fabius overview during topic selection, which is not a mathematical
dependency. It was a targeted research survey, not a full audit of every
repository file or a build of the Lean development; no repository file was
modified or uploaded. From Gao's version~1 it used \S4.5.1 (the definition of
$F_6$, its stage and sweep, determinant, and sextic fiber polynomial) and
Appendix~A.3 (the deferred fiber stratifications); it reports inspecting the
rendered formulas on PDF pages~22 and~23 and the deferred statement on
page~31. (Source~02 cites printed pages~23 and~24; the two are consistent
with an offset between PDF and printed page numbers, which this merge did
not check against the paper.) Source~05 credits the map to Gao, and claims
as new the derivative-localization identity, the exhaustive multiplicity
classification, the complete hyperplane fibers, the unique square member, the
explicit omitted surface, the stabilization obstruction from dimension three,
the global nonproperness equation and $c$-adic coefficient, the normalization
and local analytic types of the discriminant, full $S_6$ monodromy, the
intermediate-field obstruction, and its rational $S_6$ target. Its targeted
searches combined Gao, $F_6$, fiber, monodromy, omitted surface, codimension
three, and the distinctive rational coefficients of the omitted parameter; no
overlapping complete treatment was located. This limited search is not a
comprehensive priority check and does not rule out independent or newly
posted work. Source~05 credits only the construction to Gao; source~02 also
credits the degree profile and the four-point fibers over the nonzero axis.
Neither attribution was checked against Gao's paper by this merge.

\Needspace{6\baselineskip}
\section{The four sweep quotients \src{05}}
\label{f6:kf:app:polynomials}
Substituting $(w_1,w_2,w_3)=(\gamma u_1,\gamma^3u_2,\gamma^4u_3)$ into the
sweep coordinates $S_i$ of Section~\ref{f6:sec:map} and dividing by
$\gamma^i$ gives
\begin{align}
\widehat S_1={}&1-2u_1+5\gamma u_1^2+\gamma^2(-2u_1^3+2u_2)
 +\gamma^3(8u_1^4-u_3)+\gamma^4(-10u_1^5+2u_1u_3),\label{f6:kf:eq:E1}\\
\widehat S_2={}&u_1-u_1^2+\gamma(4u_1^3+u_2)+\gamma^2(-2u_1^4+2u_1u_2)
 +8\gamma^3u_1^5+\gamma^4(-10u_1^6+2u_1^2u_3),\label{f6:kf:eq:E2}\\
\widehat S_3={}&u_1^2+u_2+\gamma(3u_1^4+2u_1u_2)+\gamma^2(-2u_1^5+2u_1^2u_2)\notag\\
 &+\gamma^3(8u_1^6+u_1^2u_3)+\gamma^4(-10u_1^7+2u_1^3u_3),\label{f6:kf:eq:E3}\\
\widehat S_4={}&u_1^4+(1-2u_1)u_2+u_3+3\gamma u_1^2u_2+\gamma^2(-u_1^6+u_2^2)\notag\\
 &+\gamma^3(8u_1^4u_2-2u_1^3u_3-u_2u_3)+\gamma^4(-10u_1^5u_2+2u_1^4u_3)
 -\gamma^6u_1^2u_3^2.\label{f6:kf:eq:E4}
\end{align}
These expressions have $9$, $9$, $10$ and $13$ expanded monomials before the
stage substitution. They suffice to reproduce the entire map, the
divisibility proof of Section~\ref{f6:sec:map}, and the exceptional-hyperplane
calculation of Section~\ref{f6:kf:sec:hyperplane}, without printing the large
final coordinate expansions. (Source~05 abbreviates $(u_1,u_2,u_3)$ as
$(u,v,w)$ in this appendix only; this report does not, since $v$ and $w$ are
the chart and root variables here.)

\Needspace{6\baselineskip}
\section{Provenance of this merged report}
\label{f6:app:provenance}
This report was built in the write phase of ProveIt's incoming batch~36 from
two manuscripts of that batch, both AI-assisted (prepared with ChatGPT), both
dated 24~September~2026, and both pinned to ProveIt commit
\code{e21766d04c2b8a9b2cdba0cd43e563024b3bd1b9}. They were placed in commit
\code{1a1396d4d}.
\begin{itemize}
\item \emph{Source~02} (batch~36 manuscript~02; archive
\code{ProveIt_Six_Sheets_Research}; main file
\code{Six_Sheets_Three_Escapes.tex}, 22-page PDF): the base. It contributes
the text of Sections~\ref{f6:sec:intro}--\ref{f6:sec:conclusion} and
Appendices~\ref{f6:app:reproduce}--\ref{f6:app:sources} apart from the
tagged insertions: the uniform chart across $c=0$, the explicit finite
normalization and its boundary, the normalized discriminant and trace form,
the nonproperness locus, the general quadratic-contraction lemma, geometric
and arithmetic $S_6$ (primes $3,13,37$), the real fiber sizes, the inversion
procedure, its Lean plan and questions~1--10.
\item \emph{Source~05} (batch~36 manuscript~05; archive
\code{ProveIt_Keller_Fibers_Research}; main file \code{article.tex}, 24-page
PDF): merged in. It contributes the material tagged \src{05}: the multiplicity
patterns and the three-distinct-roots lemma, the reconstruction over $c\ne0$
and the finite completion, the hyperplane computed on the source, the
stabilization obstruction, the normalization of $\dd$, transversality along
$c=0$, the local analytic types, normal crossings and $(\Z/2)^3$ along $M$, a
second Morse slice, the intermediate-field corollary, a second arithmetic
certificate (primes $7,11,269$), its exact checks, its layered Lean plan and
directions~11--18.
\end{itemize}

\emph{Choices made in the merge.}
\begin{enumerate}[label=\textup{(\arabic*)}]
\item Source~02 is the base: its chart is a scheme isomorphism across the
whole target, its normalization lemma holds over any domain, and its finite
normalization answers source~05's question on a global completion
(direction~11).
\item Source~05 is printed in source~02's notation (Table~\ref{f6:tab:dictionary});
no normalization changed. Renamed symbols of source~05: $(C,p_1,\dots,p_4)\to
(c,A,B,D,E)$; $\xi_i\to a_i$; $(t,s,r)\to(w_1,w_2,w_3)$; $E_i\to\widehat S_i$;
$J_i\to\beta_i$ (the section coordinates); $R_{a,b,c}(T)\to P$ with root $w$ and
coefficients $(a,b,d)$; $D\to\Delta$; $\widetilde D\to H$; $H\to A^2+4B$;
$\eta\to t$; $\mathcal M\to M$; $S_F\to\NP(\F)$; $Q\to Q_\star$;
$H_4\to h_4$; $K\to\Lambda$; power sums $p_j\to N_j$; $\delta\to\theta$ (two
roots); $(u,v,w)$, $Z\to(\sigma_1,\sigma_2,\sigma_3)$, $\omega$ (local types);
$Z_i\to\omega_i$ (root clusters); $k$ (count of vanishing $\epsilon_i$) $\to j$;
$P(T)$, $c\to\hat p(w)$, $\tau$ (Morse slice); $q_*\to b_{**}$; $N\to\Omega_\Phi$;
the lifted factors $A_C,B_C\to\phi_c,\psi_c$; $d,e,f\to e_1,e_2,e_3$ (the
cubic's coefficients); $d\to d_1$ (the gcd in the fiber count);
$k[s,t]\to k[y_1,y_2]$; the
Morse-lemma degree $d\to n$; and, in Appendix~\ref{f6:kf:app:polynomials},
$(u,v,w)\to(u_1,u_2,u_3)$.
\item One symbol of source~02 was renamed: its nonproperness set
$\mathcal N_{\F}$, which clashed with its normalization algebra
$\normal=\mathcal N$, is $\NP(\F)$ here.
\item Source~02's typographical defect in \eqref{f6:eq:fixedcoeff}, where
``$\backslash$text'' had been corrupted into a tab followed by ``ext'' (its
PDF printed ``exttermsof degreeatmosttwo''), is corrected.
\item Results proved by both sources are printed once and credited to both:
polynomiality and $\det J=-290$, the section at $x=0$, the derivative
identity, the reconstruction, the gcd count, the spectrum, the unique square,
$M$, $\Delta=c^2H$ with \eqref{f6:eq:H0}, the irreducibility of $\Delta$ and
$H$, the nonproperness locus, geometric $S_6$ and the radical obstruction.
Genuinely different proofs of source~05 are kept as marked second routes:
by finite jets (polynomiality), by explicit expansion (derivative identity),
on the source (hyperplane count), through the effective parameters (omitted
surface), through Newton sums (unique square), by coprime factor lifting
($c$-adic factor), through the finite completion (nonproperness), and on a
different Morse line (geometric $S_6$). Source~05's arithmetic certificate
is a different target and is printed as its own theorem.
\item Source~05's corollary on the rational omitted point is printed as a
remark, since the point is already the zero-point witness of
Table~\ref{f6:tab:witness}. Its Lemma on connected covers is printed with
the proof of Theorem~\ref{f6:thm:geometricS6}, which contains the same
argument. Its appendix of Frobenius tests is printed in
Section~\ref{f6:kf:sec:arithmetic}.
\item Source~05's eight directions are printed as 11--18. Direction~11 (a
global finite completion) is answered by source~02 and marked so;
direction~15 (real chambers) is re-scoped after source~02's
Theorem~\ref{f6:thm:real}; source~02's question~3 is re-scoped after
source~05's local types. No other question is answered by the other source.
\item Source~02's source-audit appendix and source~05's audit are kept as
historical records of the pinned commit; the files they inspected are
unchanged since the pin. The reproducibility appendix was rewritten for the
shipped file names.
\end{enumerate}
Label prefixes: every label carries \code{f6:}. Labels of source~02 keep
their names after the prefix; labels introduced for source~05's material
carry \code{f6:kf:}; the five other labels added by the merge (four on editorial text, one on
source~02's previously unlabelled Conclusion)
(\code{f6:sec:merge}, \code{f6:sec:dictionary}, \code{f6:tab:dictionary},
\code{f6:sec:conclusion} and \code{f6:app:provenance}) carry \code{f6:} only.

\begin{thebibliography}{9}
\addcontentsline{toc}{section}{References}
\bibitem{ProveIt}
Vladimir Reshetnikov, \emph{ProveIt}, GitHub repository, Jacobian-conjecture
project and research README; revision
\code{e21766d04c2b8a9b2cdba0cd43e563024b3bd1b9}.
\url{https://github.com/VladimirReshetnikov/ProveIt}.
Source~05 cites the same revision, in particular the Lean module
\code{Counterexample.lean} and the research overview in
\code{Algebra/JacobianConjecture}.

\bibitem{ProveItWalkthrough}
Vladimir Reshetnikov, \emph{ProveIt Walkthrough}, source
\code{ProveIt_Walkthrough.tex} at the same revision; its Jacobian discussion
cites Gao's account as external context, separate from its local formal
identities (cited by source~05).
\url{https://github.com/VladimirReshetnikov/ProveIt/blob/e21766d04c2b8a9b2cdba0cd43e563024b3bd1b9/ProveIt_Walkthrough.tex}.

\bibitem{Gao}
Shuhong Gao, \emph{Counterexamples to the Jacobian conjecture in dimensions
greater than two}, arXiv:2608.00222v1, 31 July 2026.
\url{https://arxiv.org/abs/2608.00222v1}.
In particular \S4.5.1, Theorem 4.5, and Appendix A.3.

\bibitem{StacksEtale}
The Stacks Project Authors, \emph{The Stacks Project}, Tag 02GH,
``\'Etale morphisms,'' including the lemma that a morphism between
schemes \'etale over a common base is \'etale.
\url{https://stacks.math.columbia.edu/tag/02GH}.

\bibitem{StacksOpen}
The Stacks Project Authors, \emph{The Stacks Project}, Tag 025G,
Theorem 41.14.1: universally injective \'etale morphisms are open immersions.
\url{https://stacks.math.columbia.edu/tag/025G}.

\bibitem{StacksDisc}
The Stacks Project Authors, \emph{The Stacks Project}, Tag 0BJF,
Lemma 49.3.1: the discriminant of a finite locally free morphism and
its \'etale locus.
\url{https://stacks.math.columbia.edu/tag/0BJF}.

\bibitem{Milne}
James S. Milne, \emph{Fields and Galois Theory}, version 5.10,
September 2022, 144 pp. Theorem 4.28, Lemma 4.31, and Chapter 5.
\url{https://www.jmilne.org/math/}.
\end{thebibliography}
\end{document}
